% !TEX program = xelatex
% ============================================================
%  SY02-45  几何 Langlands（据 Gaitsgory、Raskin 等）
%  Peter Scholze，Séminaire BOURBAKI，78e année (2025–2026)，n° 1252，Mars 2026
%  编译：xelatex 三遍
% ============================================================
\documentclass[11pt,a4paper]{article}
% ============================================================
%  SY02-45  Scholze, Geometric Langlands [after Gaitsgory, Raskin, ...]
%  Séminaire Bourbaki n° 1252 (Mars 2026)  —— 中译本 preamble
% ============================================================
\usepackage[UTF8,fontset=windows,heading=false]{ctex}
\usepackage{amsmath,amssymb,amsthm,mathtools,bm}
\usepackage[margin=2.6cm]{geometry}
\usepackage{xcolor}
\usepackage[hidelinks]{hyperref}
\usepackage{tikz-cd}

% ---------- 版式 ----------
\linespread{1.18}
\setlength{\parskip}{0.25em}
\setlength{\marginparwidth}{2.1cm}
\renewcommand{\proofname}{证明}
\renewcommand{\qedsymbol}{$\square$}

% ---------- 编号环境：全节共用一枚计数器 ----------
% 原文体例为「类型名在前」：Theorem 2.2. — ...，同一节内所有编号条目共用序号，
% 因此这里让九类环境共享 [section] 计数器，逐条照原文顺序排版即可还原 2.1、2.2、…
\theoremstyle{definition}
\newtheorem{theorem}{定理}[section]
\newtheorem{remark}[theorem]{注}
\newtheorem{conjecture}[theorem]{猜想}
\newtheorem{definition}[theorem]{定义}
\newtheorem{proposition}[theorem]{命题}
\newtheorem{corollary}[theorem]{推论}
\newtheorem{lemma}[theorem]{引理}
\newtheorem{example}[theorem]{例}
\newtheorem{assumption}[theorem]{假设}
\newtheorem{problem}[theorem]{问题}

% ---------- 原文页码锚（页眉 1252–NN，与 PDF 页一致）----------
\newcommand{\srcpage}[1]{\marginpar{\raggedright\footnotesize\ttfamily\color{gray}1252--#1}}

% ---------- 译注：\yizhu{内容}{等级（出处）} ----------
% 等级只许三类：术语背景 / 原文省略 / 原文笔误；出处只许 A·B·I 或「上下文推断」
\newcommand{\yizhu}[2]{\footnote{\textsf{【译注·#2】}#1}}

% ---------- 常用对象 ----------
\newcommand{\Bun}{\operatorname{Bun}}
\newcommand{\Hck}{\operatorname{Hck}}
\newcommand{\Gr}{\operatorname{Gr}}
\newcommand{\LocSys}{\operatorname{LocSys}}
\newcommand{\Rep}{\operatorname{Rep}}
\newcommand{\IndCoh}{\operatorname{IndCoh}}
\newcommand{\Coh}{\operatorname{Coh}}
\newcommand{\Perf}{\operatorname{Perf}}
\newcommand{\Perv}{\operatorname{Perv}}
\newcommand{\Sch}{\mathrm{Sch}}
\newcommand{\Corr}{\mathrm{Corr}}
\newcommand{\Fun}{\operatorname{Fun}}
\newcommand{\End}{\operatorname{End}}
\newcommand{\Hom}{\operatorname{Hom}}
\newcommand{\Map}{\operatorname{Map}}
\newcommand{\Spec}{\operatorname{Spec}}
\newcommand{\Pic}{\operatorname{Pic}}
\newcommand{\Sht}{\operatorname{Sht}}
\newcommand{\KL}{\operatorname{KL}}
\newcommand{\Op}{\operatorname{Op}}
\newcommand{\Whit}{\operatorname{Whit}}
\newcommand{\Sat}{\operatorname{Sat}}
\newcommand{\Frob}{\operatorname{Frob}}
\newcommand{\tr}{\operatorname{tr}}
\newcommand{\Eis}{\operatorname{Eis}}
\newcommand{\CT}{\operatorname{CT}}
\newcommand{\Higgs}{\operatorname{Higgs}}
\newcommand{\Aut}{\operatorname{Aut}}
\newcommand{\id}{\mathrm{id}}
\newcommand{\pr}{\mathrm{pr}}
\newcommand{\crit}{\mathrm{crit}}

% 对偶群与对偶环面
\newcommand{\Gd}{\check{G}}
\newcommand{\Bd}{\check{B}}
\newcommand{\Td}{\check{T}}
\newcommand{\Ud}{\check{U}}
\newcommand{\Md}{\check{M}}

% 圈群
\newcommand{\LG}{\mathrm{L}G}
\newcommand{\LpG}{\mathrm{L}^{+}G}
\newcommand{\LU}{\mathrm{L}U}
\newcommand{\LT}{\mathrm{L}T}
\newcommand{\LB}{\mathrm{L}B}
\newcommand{\LUd}{\mathrm{L}\check{U}}

% 花体
\newcommand{\cA}{\mathcal{A}}
\newcommand{\cB}{\mathcal{B}}
\newcommand{\cC}{\mathcal{C}}
\newcommand{\cD}{\mathcal{D}}
\newcommand{\cE}{\mathcal{E}}
\newcommand{\cF}{\mathcal{F}}
\newcommand{\cG}{\mathcal{G}}
\newcommand{\cH}{\mathcal{H}}
\newcommand{\cI}{\mathcal{I}}
\newcommand{\cJ}{\mathcal{J}}
\newcommand{\cK}{\mathcal{K}}
\newcommand{\cL}{\mathcal{L}}
\newcommand{\cM}{\mathcal{M}}
\newcommand{\cN}{\mathcal{N}}
\newcommand{\cO}{\mathcal{O}}
\newcommand{\cP}{\mathcal{P}}
\newcommand{\cQ}{\mathcal{Q}}
\newcommand{\cR}{\mathcal{R}}
\newcommand{\cS}{\mathcal{S}}
\newcommand{\cT}{\mathcal{T}}
\newcommand{\cU}{\mathcal{U}}
\newcommand{\cV}{\mathcal{V}}
\newcommand{\cW}{\mathcal{W}}
\newcommand{\cX}{\mathcal{X}}
\newcommand{\cY}{\mathcal{Y}}
\newcommand{\cZ}{\mathcal{Z}}

% 哥特体
\newcommand{\gc}{\mathfrak{g}}
\newcommand{\bc}{\mathfrak{b}}
\newcommand{\uc}{\mathfrak{u}}
\newcommand{\nc}{\mathfrak{n}}
\newcommand{\tc}{\mathfrak{t}}
\newcommand{\pc}{\mathfrak{p}}
\newcommand{\mc}{\mathfrak{m}}

% 数系与域
\newcommand{\Z}{\mathbb{Z}}
\newcommand{\Q}{\mathbb{Q}}
\newcommand{\R}{\mathbb{R}}
\newcommand{\C}{\mathbb{C}}
\newcommand{\F}{\mathbb{F}}
\newcommand{\A}{\mathbb{A}}
\newcommand{\Gm}{\mathbb{G}_m}
\newcommand{\Ql}{\overline{\mathbb{Q}}_{\ell}}
\newcommand{\Zl}{\mathbb{Z}_{\ell}}
\newcommand{\Fq}{\mathbb{F}_{q}}
\newcommand{\Fp}{\mathbb{F}_{p}}

% 形变与下标
\newcommand{\dR}{\mathrm{dR}}
\newcommand{\et}{\mathrm{et}}
\newcommand{\Betti}{\mathrm{Betti}}
\newcommand{\qc}{\mathrm{qc}}
\newcommand{\Nilp}{\mathrm{Nilp}}
\newcommand{\Ran}{\mathrm{Ran}}
\newcommand{\Weil}{\mathrm{Weil}}
\newcommand{\coarse}{\mathrm{coarse}}
\newcommand{\restr}{\mathrm{restr}}
\newcommand{\irred}{\mathrm{irred}}
\newcommand{\monfree}{\mathrm{monfree}}
\newcommand{\lisse}{\mathrm{lisse}}
\newcommand{\prim}{\mathrm{prim}}
\newcommand{\suave}{\mathrm{suave}}
\newcommand{\Lisse}{\operatorname{Lisse}}
\newcommand{\Ind}{\operatorname{Ind}}
\newcommand{\PrL}{\mathrm{Pr}^{L}}
\newcommand{\e}{\mathit{e}}
\newcommand{\Th}{\Theta}
\newcommand{\Om}{\Omega}
\newcommand{\OmX}{\Omega^{1}_{X}}
\newcommand{\Dmod}{D\text{-}\mathrm{mod}}
\newcommand{\perf}{\mathrm{perf}}
\newcommand{\pt}{\mathrm{pt}}
\newcommand{\loc}{\mathrm{loc}}
\newcommand{\alg}{\mathrm{alg}}
\newcommand{\der}{\mathrm{der}}
\newcommand{\op}{\mathrm{op}}
\newcommand{\ee}[1]{\mathrm{e}^{#1}}

% ---------- 排版小工具 ----------
% 行内公式中不宜断行的长串
\newcommand{\dsp}{\displaystyle}
% 参考文献条目
\newcommand{\refitem}[2]{\par\noindent\hangindent=2em\hangafter=1 #1\ #2\par}


\title{几何 Langlands\\[0.4em]（据 Gaitsgory、Raskin 等）}
\author{Peter Scholze\\[0.3em]
\small Séminaire BOURBAKI，第 78 年（2025--2026），第 1252 期，2026 年 3 月}
\date{}

\begin{document}
\maketitle

\noindent\small
\textbf{底本}：Peter Scholze, \emph{Geometric Langlands [after Gaitsgory, Raskin, \dots]},
Séminaire BOURBAKI, 78e année (2025--2026), n° 1252, Mars 2026, 56 页。
pdfTeX 排版，页眉编号 \texttt{1252--NN} 与 PDF 页码一致，译文边注中的页码即指此页。

\medskip
\noindent\small
\textbf{译注约定}：原文含糊、省略或笔误之处，以脚注标出，脚注以「【译注·类别】」起首，
类别只取\emph{术语背景}、\emph{原文省略}、\emph{原文笔误}三种，并注明依据等级
（A：权威文献；B：nLab/Stacks/ProofWiki 一类百科；I：译者按上下文推断）。
译文不补写原文所无的论证。

\medskip
\noindent\small
\textbf{编号}：原文的定理、命题、引理、推论、猜想、定义、注、例、假设、问题各成条目，
同一节内共用序号，译文沿用其编号（如「定理 2.11」「注 3.1」），未作改动。

\normalsize
\tableofcontents
\newpage

\section{引言}\srcpage{1}

Langlands 纲领是一张影响深远的猜想与定理之网，它在数学的许多不同领域中都有化身。虽然它最初由 Langlands 提出，用以处理数论中的问题，例如推广类域论的非交换互反律的存在性，以及算术局部对称空间上的自守形式理论。但此后它又在许多不同的情形中找到了化身。这些翻译通常依照 Weil 的「Rosetta Stone」（罗塞塔石碑）进行，它有三栏：
\begin{enumerate}
\item 数域的原初情形；
\item 有限域 $\Fq$ 上函数域的、与之类似而更易接近的情形；等价地，$\Fq$ 上光滑射影曲线 $X$ 的情形；
\item 代数闭域上光滑射影曲线的几何情形；例如在 $\C$ 上，这些等价于 Riemann 曲面。
\end{enumerate}

\srcpage{2}
从 1970 年代 Drinfeld 的工作开始，人们研究了 $\Fq$ 上以及其代数闭包 $\overline{\F}_q$ 上光滑射影曲线 $X$ 的 Langlands 纲领。此后数年间，人们也研究了 $\C$ 上光滑射影曲线的类似结构；但有一段时间，这些研究似乎离最初关心的数论情形越来越远，同时又发现了与数学物理的新联系。当我还是研究生、在算术 Langlands 纲领中工作时，我无法理解几何 Langlands 纲领所研究的结构究竟是什么，为什么人们应当对它们感兴趣，以及它们与最初关心的问题有何关联。

然而，在过去几年里，情况发生了剧变：（特征 $0$、处处非分歧、整体）几何 Langlands 猜想已被证明。并且，这一几何猜想又被联系回 Weil 罗塞塔石碑的第二栏，即有限域上的函数域，从而对函数域上 Langlands 纲领中一些最基本的问题给出非常精确的回答。

在本引言的余下部分，我将简要回顾 Langlands 最初在数域情形提出的问题，然后陈述 Gaitsgory–Raskin（2025a）证明的一个定理。该定理依赖许多其他论文，作者包括 Dima Arinkin、Dario Beraldo、Justin Campbell、Lin Chen、Joakim Færgeman、Dennis Gaitsgory、David Kazhdan、Kevin Lin、Sam Raskin、Nick Rozenblyum、Yakov Varshavsky——它对该函数域情形中的类似问题给出了非常精确的回答。

就其核心而言，Langlands 纲领试图描述自守函数构成的向量空间。这不过是对某个极为基本之物给出的一个花哨名称。也就是说，数学中出现的一些最基本、最突出、最漂亮的空间例子，正是那些局部以某个给定几何（例如双曲几何或球面几何）为模型的空间。更确切地说，人们固定一个 Riemann 对称空间 $M$\footnote{习惯上把该对称空间记为 $X$，但这与几何 Langlands 中所用曲线的标准记号冲突。}及其自同构李群 $G$，然后考虑局部同构于 $M$ 的流形 $N$，其坐标变换属于 $G$。若 $N$ 连通且完备，则它可以写成局部对称空间 $M/\Gamma$，其中 $\Gamma=\pi_1(N)$ 通过一个嵌入作为 $G$ 的离散子群作用。著名的刚性定理表明，这一情形实际上在大多数情况下比它看起来更具代数与数论色彩——也就是说，$G$ 带有约化线性代数群的结构，而 $\Gamma\subset G$ 是算术群。

除了一些细节，自守函数就是算术群 $\Gamma$ 对应的局部对称空间 $M/\Gamma$ 上的函数。在本引言中，我将忽略所考虑的函数空间究竟为何这一问题，而直接把 $M/\Gamma$ 上某个合适的性质良好的函数构成的向量空间记为 $\cA(M/\Gamma)$。注意，李群 $G$ 作用在 $M$ 上；无穷小作用特别地通过微分给出李代数 $\gc$ 在 $\cA(M)$ 上的作用。然而，这一作用不与 $\Gamma$ 交换，因此不在 $\cA(M/\Gamma)$ 上诱导作用。但我们可以把 $\gc$-作用延拓到泛包络代数 $U(\gc)$ 上。由

\srcpage{3}
Harish-Chandra 的工作，$U(\gc)$ 有很大的中心 $Z(U(\gc))$，它同构于 Cartan 子代数 $\mathfrak{h}\subset\gc$ 的泛包络代数中的 Weyl 群不变量 $U(\mathfrak{h})^W$。由于它的中心性，$Z(U(\gc))$ 在 $\cA(M)$ 上的作用现在与 $G$ 的作用、从而与 $\Gamma$ 的作用交换，于是它在 $\cA(M/\Gamma)$ 上诱导作用。特别地，这给出 Laplace 算子的作用。

另一方面，由于 $\Gamma$ 是算术的，它与许多共轭子群可公度。为简化讨论，让我作如下假设。

\begin{assumption}
群 $G$ 分裂且单连通。从现在起，我将用 $G$ 记它在 $\Z$ 上的分裂形式，于是先前的 $G$ 现在称为 $G(\R)$。此外，$\Gamma=G(\Z)$ 是 $G$ 的全部整点构成的群。
\end{assumption}

对任意 $g\in G(\Q)$，群 $\Gamma_g=\Gamma\cap g^{-1}\Gamma g$ 是 $\Gamma$ 的有限指数子群。此时，$g$ 在 $M$ 上的作用下降为一个映射
\[ M/\Gamma_g\xrightarrow{g}M/\Gamma, \]
连同自然投影 $M/\Gamma_g\to M/\Gamma$，我们得到 Hecke 对应
\[ \begin{tikzcd} & M/\Gamma_g \arrow[dl] \arrow[dr, "g"] & \\ M/\Gamma & & M/\Gamma. \end{tikzcd} \]
这诱导出自同态
\[ T_g\colon\cA(M/\Gamma)\to\cA(M/\Gamma), \]
其做法是先把函数拉回到 $M/\Gamma_g$，再沿有限映射 $g\colon M/\Gamma_g\to M/\Gamma$ 的纤维求和。这些 Hecke 算子与上面构造的微分算子交换，而且有些出人意料的是，它们彼此也交换\footnote{这里用到 $\Gamma=G(\Z)$；一般情形略为复杂。}。合起来，它们给出一个（无限生成的）交换 $\C$-代数 $\mathbb{T}=\mathbb{T}_G$ 在 $\cA(M/\Gamma)$ 上的作用。

下述问题可以说是 Langlands 纲领的中心问题。

\begin{problem}
描述带有 $\mathbb{T}$ 作用的向量空间 $\cA(M/\Gamma)$。
\end{problem}

Langlands 的进路是把 $\cA(M/\Gamma)$ 作为 $\mathbb{T}$-模作谱分解。他证明它可分解为「连续谱」与「离散谱」。此外，连续谱允许用所谓的 Eisenstein 级数给出一个（复杂但本质上是显式的）描述，这些级数由与 $G$ 的真 Levi 子群相伴的局部对称空间上的自守函数归纳地构造。「离散谱」允许分解为不可约「自守表示」的直和。其中一些仍间接地 \srcpage{4}来自 Eisenstein 级数，作为「Eisenstein 级数的留数」，而另一些则是「对 $G$ 真正新的」；这些被称为「尖点的」。

关于可能的「Hecke 特征值系统」，即 $\cA(M/\Gamma)$ 的支撑中的同态 $\psi\colon\mathbb{T}\to\C$，存在若干深刻的预测。为了陈述它们，多谈一点 Hecke 算子的结构会有用。

我们关心空间 $M/G(\Z)$，其中 $M=G(\R)/K_\infty$ 是 $G(\R)$ 的对称空间；这里 $K_\infty\subset G(\R)$ 是极大紧子群。给出该商的一个点，等价于给出一个 $G(\Z)$-挠子——等价地，给出 $\Spec(\Z)$ 上的一个 $G$-挠子 $E$——连同相伴的 $G(\R)$-挠子到 $K_\infty$ 的一个约化——换言之，$E\times_{\Spec(\Z)}\Spec(\R)$ 上的一个「度量」。例如，若 $G=\mathrm{SL}_n$，从而 $K_\infty=\mathrm{SO}_n(\R)\subset\mathrm{SL}_n(\R)$，则这是如下数据：秩为 $n$ 的 $\Z$-格 $E$，连同平凡化的行列式 $\det(E)\cong\Z$，以及 $E\otimes_\Z\R$ 上的一个正定配对。

在这个图像中，由 $g\in G(\Q)$ 标号的 Hecke 算子一般地不改变 $G$-挠子 $E$，即在 $\Spec(\Q)$ 上不改变它，而是在 $\Spec(\Z)$ 的某些有限位处修改它。实际上，这可以一次只对一个素数 $p$ 进行，使用 $g\in G(\Z[\frac1p])$ 的 Hecke 算子。$p$ 处的 Hecke 算子给出代数
\[ \C\big[G(\Z)\backslash G(\Z[\tfrac1p])/G(\Z)\big] \]
的作用，该代数由 $G(\Z[\frac1p])$ 中 $G(\Z)$-双陪集上的紧支撑函数组成，其中 $T_g$ 对应双陪集 $G(\Z)gG(\Z)$ 的特征函数。实际上，由强逼近，该代数可以用 $p$ 进群 $G(\Q_p)$ 改写成
\[ \C\big[G(\Z_p)\backslash G(\Q_p)/G(\Z_p)\big]. \]

Harish-Chandra 对泛包络代数中心的描述，其类似物是 Satake 同构：固定分裂极大环面 $T\subset G$ 及其 Weyl 群 $W$，则有 $\C$-代数同构
\[ \C\big[G(\Z_p)\backslash G(\Q_p)/G(\Z_p)\big]\cong\C[X_*(T)]^W; \]
特别地，左端交换。人们常常形式地引入 $\C$ 上的对偶环面 $\Td$，其特征群 $X^*(\Td)=X_*(T)$ 与 $T$ 的余特征群等同；于是
\[ \C[X_*(T)]^W=\cO(\Td/W), \]
即仿射簇 $\Td/W$ 上正则函数构成的代数。若再引入 Langlands 对偶群 $\Gd$，其根数据与 $G$ 的根数据对偶，特别地它以 $\Td$ 为极大环面、以 $\check W=W$ 为 Weyl 群，则我们可以进一步把它重新解释为
\[ \cO(\Gd/_{\mathrm{ad}}\Gd), \]
即 $\Gd$ 上的共轭不变函数。因此，Satake 同构给出同构
\[ \C\big[G(\Z_p)\backslash G(\Q_p)/G(\Z_p)\big]\cong\cO(\Gd/_{\mathrm{ad}}\Gd), \]
这是群 $G$ 与其 Langlands 对偶群 $\Gd$ 之间在函数层面的第一座桥梁。

\srcpage{5}
上述讨论给出抽象 Hecke 代数
\[ \mathbb{T}=\bigotimes_{p}\C\big[G(\Z_p)\backslash G(\Q_p)/G(\Z_p)\big]\otimes Z(U(\gc)) \]
在 $\cA(M/\Gamma)$ 上的作用。考虑 $\cA(M/\Gamma)$ 的支撑中的一个映射 $\psi\colon\mathbb{T}\to\C$。映射 $\psi$ 对应于对所有 $p$ 的映射
\[ \C\big[G(\Z_p)\backslash G(\Q_p)/G(\Z_p)\big]\to\C, \]
以及映射 $Z(U(\gc))\to\C$。由 Satake 同构与 Harish-Chandra 同构，它们分别给出：对每个素数 $p$，一个 $\Gd(\C)$ 中元素的 $\Gd$-共轭类 $\psi_p$；以及 $\psi_\infty\in\check{\gc}$。

有两个核心猜想；通过归纳，可以假设 $\psi$ 属于离散谱。在陈述下述猜想时，我先假设 $G=\mathrm{SL}_n$（从而 $\Gd=\mathrm{PGL}_n$），且 $\psi$ 是尖点的。
\begin{enumerate}
\item （Ramanujan）对每个 $p<\infty$，$\psi_p$ 落在 $\Gd(\C)$ 的一个紧子群中；而 $\psi_\infty$ 落在 $\Gd(\C)$ 的某个紧子群的李代数中。
\item （互反）$\Spec(\Z)$ 上存在不可约 $\Gd$-局部系统\footnote{警告：这一概念尚不知道，参看下面的注 1.4。}，其中 $\psi_p$ 是沿「圆」$\Spec(\Fp)\subset\Spec(\Z)$ 的单值性，而 $\psi_\infty$ 给出广义 Hodge 数。
\end{enumerate}
对完整的离散谱，以及对 $\mathrm{SL}_n$ 以外的群，这两个猜想都需要稍作修改；Arthur（1989）陈述了相应的推广。

\begin{remark}
Ramanujan 猜想是一个极其简单而一般的陈述。应用于 $G=\mathrm{SL}_2$，并且只看微分算子的作用，它预测在 $\mathbb{H}/\mathrm{SL}_2(\Z)$ 上 Laplace 算子的谱隙至少为 $\frac14$；实际上，它对所有算术双曲曲面都预测同样的事。此外，它把这一现象适当地推广到所有算术局部对称空间。对 $p<\infty$ 的 $\psi_p$，它预测 Hecke 特征值的「平方根消去」；在某些情形，Deligne 通过证明一个互反的版本与 Weil 猜想，即「有限域上的 Riemann 假设」，著名地证明了这一点——特别是在 Ramanujan 最初提出其猜想的权 $12$ 模形式 $\Delta$ 的情形。

注意，Ramanujan 猜想非常强烈地提示存在某个与 $\psi$ 相伴的 $\Gd(\C)$ 的典范紧子群，在共轭意义下所有 $\psi_p$ 都落在其中（而 $\psi_\infty$ 落在其李代数中）。实际上，人们预期这样的子群存在，且 $\psi_p$ 在其中等分布——这就是 Sato–Tate 猜想。此外，这个紧子群应当恰好是互反猜想所预测的 $\Gd$-局部系统的单值群。
\end{remark}

\begin{remark}
遗憾的是，互反猜想就没有同样的简单性：它只是一个元猜想。也就是说，相关的 $\Gd$-局部系统概念尚不知道。定义这样一个概念，本质上等价于定义 $\Spec(\Z)$ 上 $\C$-局部系统的某个 Tannakian 范畴，而这又等价于定义

\srcpage{6}
某个投射代数（或投射紧）群——这就是猜想中的 Langlands 群。然而，我相信人们不应当把注意力集中在 Langlands 群上，而应当寻找 $\Spec(\Z)$ 上 $\C$-局部系统的一个更一般的理论——或者说，真正要对 $\Spec(\Z)$ 上的概形建立「扭量结构的变体」的完整六函子形式化。这样的理论应当允许 Hodge 数是复数而非整数——在 $\Spec(\R)$ 上，这在 Scholze（2024）发展的理论中是可行的。

若 $\psi_\infty\in\check{\gc}$ 的共轭类在 $\check{\tc}$ 中有整的代表，从而所预测的 Hodge 数是整数，则可以使用一个非猜想的替代物：也就是说，固定某个素数 $\ell$ 与同构 $\C\cong\overline{\Q}_\ell$，并使用 étale $\overline{\Q}_\ell$-局部系统。回忆一下，$\ell$ 进局部系统本质上是 $\ell$ 进 Galois 表示，这就给出自守形式与 Galois 表示之间著名的猜想性关系\footnote{此处我忽略了一个重要的 $\rho$-shift（$\rho$-平移），关于这一点的详细阐述参看 Buzzard–Gee（2014）。}。
\end{remark}

上述猜想有一系列加强版本，它们不仅预测某个 Hecke 特征值系统 $\psi$ 的支撑，还预测它出现的重数。特别地，这预测在某种意义上 $\Spec(\Z)$ 上所有 $\Gd$-局部系统都应当来自某个自守形式。

下面的猜想实际上会对问题 1.2 给出更精确的回答。如同 Langlands 纲领中许多最深刻的猜想一样，它其实是一个猜想的模板——许多术语仍未定义\footnote{此外，还缺少一些精确化；尤其是，阿基米德位处的处理可能需要一些修正，而且我忽略了由各种 $\rho$-shift 引起的可能扭转。}。

\begin{conjecture}
存在 $\C$ 上的一个叠 $\LocSys_{\Gd}$，它参数化 $\Spec(\Z)$ 上的 $\Gd$-局部系统，并带有一个映射 $\LocSys_{\Gd}\to\Spec(\mathbb{T})$，该映射由把 $\Gd$-局部系统限制到所有 $\Spec(\Fp)\subset\Spec(\Z)$、以及广义 Hodge 数诱导。此外，存在自守函数空间与 $\LocSys_{\Gd}$ 上对偶复形 $\omega_{\LocSys_{\Gd}}$ 的整体截面之间的典范 $\mathbb{T}$-等变同构
\[ \cA(M/\Gamma)\cong\Gamma(\LocSys_{\Gd},\omega_{\LocSys_{\Gd}}). \]
\end{conjecture}

把这个猜想拆成两部分是有益的：
\begin{enumerate}
\item 可以把 $\cA(M/\Gamma)$ 提升为 $\LocSys_{\Gd}$ 上的拟凝聚层。
\item 这个拟凝聚层就是对偶复形。
\end{enumerate}
第一部分已经给出 Langlands 互反的「自守到 Galois」方向，它蕴含 $\cA(M/\Gamma)$ 的支撑必须包含在 $\LocSys_{\Gd}\to\Spec(\mathbb{T})$ 的像中。第二部分则相当于描述给定的 $\Gd$-局部系统出现的重数，并包含「Galois 到自守」的反方向。

\begin{remark}
该猜想的以下方面尚不清楚：

\srcpage{7}
\begin{enumerate}
\item $\Gd$-局部系统的概念是什么？
\item $\LocSys_{\Gd}$ 是哪种「叠」？
\end{enumerate}
我们注意到，这里肯定需要某种解析几何，这已经是因为：所寻求的向量空间 $\cA(M/\Gamma)$ 带有非平凡拓扑，因此必须从泛函分析的观点来研究。
\end{remark}

受下文将讨论的几何 Langlands 联系的启发，近年来若干作者群体提出了类似的猜想，例如 Zhu（2025，第 4.7 节）与 Emerton。不过，通常不是针对通常的自守形式，而是针对 $p$ 进自守形式（例如以 $M/\Gamma$ 的 $p$ 进奇异上同调的形式出现）来表述。在那种情形，所需的 $\Gd$-局部系统模空间可以（本质上）\footnote{即便在那里，阿基米德位也可能引起一些问题；不过，可以用 Clausen 最近的 Weil–Moore anima 理论来修正。}用 $p$ 进 étale 局部系统来定义。

猜想 1.5 会比 Arthur–Langlands 的猜想更精确：也就是说，它用 $\Gd$-局部系统描述带有 $\mathbb{T}$-作用的实际的自守形式向量空间；而 Langlands 只着眼于描述出现的自守表示的同构类与重数。令人惊讶的是，Arthur 精细化的参数概念并不（也不应当）进入它的陈述。相反，当人们分析右端的离散谱究竟是什么时，明显的部分来自 $\LocSys_{\Gd}$ 的孤立点（这些对应不可约 $\Gd$-局部系统），但进一步的贡献来自对偶复形 $\omega$ 变得更复杂的奇点。对这一现象的分析应当能恢复 Arthur 对 Langlands 猜想的精细化\footnote{正如我将在下文讨论的，在函数域情形这会成为一个精确的问题，它在 Gaitsgory–V. Lafforgue–Raskin 最近的工作中得到处理，参看 Raskin（2025）。}。

Gaitsgory、Raskin 等人的工作处于这一情形的函数域类比之中。粗略地说，这相当于把 $\Z$ 换成 $\Fp[\mathbf{t}]$，把 $\R$ 换成 $\Fp((\mathbf{t}^{-1}))$。注意，正如 $\Z\subset\R$ 是离散余紧子群，$\Fp[\mathbf{t}]\subset\Fp((\mathbf{t}^{-1}))$ 也是离散余紧子群。$M=G(\R)/K_\infty$ 的类比是 $G(\Fp((\mathbf{t}^{-1})))/G(\Fp[[\mathbf{t}^{-1}]])$，其中 $G(\Fp((\mathbf{t}^{-1})))$ 是以 $G(\Fp[[\mathbf{t}^{-1}]])$ 为极大紧子群的局部紧群。$\Gamma=G(\Z)$ 的类比是 $G(\Fp[\mathbf{t}])$。我们于是被引向考察双商
\[ G(\Fp[\mathbf{t}])\backslash G(\Fp((\mathbf{t}^{-1})))/G(\Fp[[\mathbf{t}]]). \]
结果表明它可以有更简洁的描述：它就是 $\mathbb{P}^1_{\Fp}$ 上 $G$-挠子的同构类集合。事实上，这不过是在镜像我们先前用 $\Spec(\Z)$ 上的 $G$-挠子连同 $\Spec(\R)$ 上的一个「度量」对 $M/\Gamma$ 的描述。

正如 $\Z$ 上的 Langlands 纲领对任意数域 $K$ 有关于 $\cO_K$ 的变体，我们在函数域情形也可以允许 $\mathbb{P}^1_{\Fp}$ 的任意有限覆盖，即对 $\Fp$ 上任意光滑射影曲线 $X$ 工作。设 $\Fq$ 为其常数域，即 $\Fp$ 的有限扩张。由 Hecke 的一个定理的函数域类比（参看 Weil（1974，第 291 页）），人们总能找到「Theta 特征」，即 $X$ 上线丛 $\OmX$ 的一个平方根 $\Theta$，我们将其固定\footnote{这一点只在稍后的归一化中用到，参看注 1.8。}。

\srcpage{8}
自守函数空间的类比是函数空间 $\cA(\Bun_G(\Fq))$，即 $X$ 上 $G$-丛的模空间 $\Bun_G$ 的 $\Fq$-点上的函数空间；从第 2 节开始将更详细地讨论它。Langlands 纲领的几乎所有方面都忠实地翻译到函数域情形，但许多事情容易得多：例如，$\Bun_G(\Fq)$ 只是一个离散点集，因此作为抽象向量空间，$\cA(\Bun_G(\Fq))$ 的描述是平凡的——这已经去掉了所有非平凡的分析。但同样地，有 Hecke 算子作用在 $\cA(\Bun_G(\Fq))$ 上。更确切地说，对每个点 $x\in X$，我们得到 Hecke 对应
\[ \begin{tikzcd} & \Hck_{G,x}(\Fq) \arrow[dl, "p_1"'] \arrow[dr, "p_2"] & \\ \Bun_G(\Fq) & & \Bun_G(\Fq), \end{tikzcd} \]
其中 $\Hck_{G,x}$ 参数化 $G$-丛对 $E,E'$，连同它们在 $X\setminus\{x\}$ 上的一个同构。这在 $\cA(\Bun_G(\Fq))$ 上诱导出 Hecke 代数
\[ \mathbb{T}_x=\C\big[G(\cO_x)\backslash G(F_x)/G(\cO_x)\big]\cong\cO(\Gd/_{\mathrm{ad}}\Gd) \]
的作用。这里 $\cO_x$ 表示 $X$ 在 $x$ 处的完备局部环——一个完备离散赋值环——而 $F_x$ 是它的分式域；这完全类似于 $\Z_p\subset\Q_p$，其中 $x$ 扮演 $p$ 的角色。也就是说，$h\in\C[G(\cO_x)\backslash G(F_x)/G(\cO_x)]$ 通过如下公式作用：
\[ f\in\cA(\Bun_G(\Fq))\mapsto\Big(x\mapsto\sum_{y,\ p_2(y)=x}h(y)f(p_1(y))\Big)\in\cA(\Bun_G(\Fq)). \]

同样地，Arthur–Langlands 的猜想预测用 $X$ 上 $\Gd$-局部系统来描述。仍然没有 $\C$-局部系统的概念，但在函数域情形，人们可以直接使用 étale $\overline{\Q}_\ell$-局部系统的概念，任选一个同构 $\C\cong\overline{\Q}_\ell$——根本原因在于，在函数域情形，自守函数的概念是纯代数的。

Zhu（2025）的工作表明，存在一个性态良好的（导出 Artin）叠 $\LocSys_{\Gd}$，它定义在 $\overline{\Q}_\ell$ 上，参数化 $\Gd$-局部系统，并带有一个到 $\Spec(\mathbb{T})$ 的映射；这里 $\mathbb{T}=\bigotimes_{x\in X}\mathbb{T}_x$ 是抽象 Hecke 代数。下述定理（本质上）在函数域情形解决了猜想 1.5：

\begin{theorem}[Gaitsgory 与 Raskin，2025a，定理 0.2.2]
设 $G$ 是 $\Fq$ 上的约化群，并假设 $p$ 相对于 $G$ 充分大。则存在既开又闭的子叠 $\LocSys'_{\Gd}\subset\LocSys_{\Gd}$，以及典范的 $\mathbb{T}$-等变同构
\[ \cA_c(\Bun_G(\Fq))\cong\Gamma(\LocSys'_{\Gd},\omega_{\LocSys'_{\Gd}}). \]
\end{theorem}

\srcpage{9}
这里 $\cA_c$ 表示紧支撑函数。Raskin（2025）宣布了与 Gaitsgory、V. Lafforgue 的合作工作，在函数域情形推出 Arthur–Ramanujan 猜想。因此，归根结底，人们证明的是关于实际数字的具体陈述，即 Hecke 算子的特征值！

当然，人们预期可以取 $\LocSys'_{\Gd}=\LocSys_{\Gd}$；这对 $G=\mathrm{SL}_n$ 是已知的。

\begin{remark}[Whittaker 归一化]
这个定理有一个令人惊讶的方面：它给出向量空间之间一个特别的同构。归一化的一个关键方面是下述匹配，它用到我们对 Theta 特征 $\Theta$ 的选取，以及在 $\C$ 中单位根的标准选取。在右端，$\LocSys_{\Gd}$ 的几何性质（其余切复形的行列式典范平凡）保证 $\omega_{\LocSys_{\Gd}}$ 带有一个典范截面，于是右端有一个特别的向量。Gaitsgory–V. Lafforgue–Raskin 证明，这对应于 Poincaré 级数
\[ \mathcal{P}\in\cA_c(\Bun_G(\Fq)), \]
它是把 Whittaker 特征
\[ \Bun_U^{\Th}(\Fq)\to\C^\times \]
沿 $\Bun_U^{\Th}(\Fq)\to\Bun_G(\Fq)$ 的纤维求和而得到的。这里 $\Bun_U^{\Th}$ 是 $U$-丛模空间的一个扭转形式，其中 $U$ 是 Borel 子群 $B\subset G$ 的单幂根。更确切地说，$\Bun_U^{\Th}$ 是 $U^{\Th}$-丛的模空间，这里 $U^{\Th}$ 是 $X$ 上的一个群概形，即 $U$ 被 $T$-挠子 $2\rho(\Theta)$ 扭转所得的扭转形式，而该挠子是线丛 $\Theta$ 所对应的 $\Gm$-挠子在映射 $2\rho\colon\Gm\to T$ 下的前推。这一扭转保证：对任意单根 $\alpha$，$U^{\Th}$ 的相应根空间不是 $\mathbb{G}_a$ 而是 $\OmX$，于是我们得到一个投影
\[ p_\alpha\colon U^{\Th}\to\OmX. \]
把所有 $p_\alpha$ 相加，我们得到投影 $U^{\Th}\to\OmX$。过渡到分类空间并从 $X$ 取映射叠，这给出映射
\[ \Bun_U^{\Th}(\Fq)\to H^1(X,\OmX), \]
而 Serre 对偶给出同构 $H^1(X,\OmX)=\Fq$。最后，Whittaker 特征是复合
\[ \Bun_U^{\Th}(\Fq)\to\Fq\xrightarrow{\tr}\Fp\hookrightarrow\C^\times, \]
其中用到 $\C$ 中单位根的标准选取。
\end{remark}

如上所述，该定理是数十年来几何 Langlands 纲领研究的成果。几何 Langlands 纲领的起点是 Drinfeld 的观察：在某些情形，可以取一个 $\Gd$-局部系统，并在另一侧产生相应的自守形式，参看 Drinfeld（1983）。

\srcpage{10}
为此，Drinfeld 实际上构造了 $\Bun_G$ 上的一个 $\ell$ 进层 $\cF$，然后通过 Grothendieck 的层-函数字典得到 $\Bun_G(\Fq)$ 上的函数 $\tr_\cF$：对任意点 $x\in\Bun_G(\Fq)$，对应 $\tr(\Frob_x|\cF_{\bar x})$，即（$x$ 上方一个几何点 $\bar x$ 处的）Frobenius 迹。这提示，在几何上，$X_{\overline{\F}_q}$ 上的 $\Gd$-局部系统与 $\Bun_{G,\overline{\F}_q}$ 上的可构造层之间可能存在关系。

非常粗略地说，定理 1.7 的证明现在如下：
\begin{enumerate}
\item 把几何 Langlands 对应表述为范畴等价，大致说它把 $\Bun_{G,\overline{\F}_q}$ 上的 étale $\ell$ 进层
\[ D_{\et}(\Bun_{G,\overline{\F}_q},\overline{\Q}_\ell) \]
与 $\Gd$-局部系统的叠 $\LocSys^{\mathrm{geom}}_{\Gd}$ 上的拟凝聚层
\[ D_{\qc}(\LocSys^{\mathrm{geom}}_{\Gd}) \]
联系起来，其中 $\LocSys^{\mathrm{geom}}_{\Gd}$ 是 $X_{\overline{\F}_q}$ 上 $\Gd$-局部系统的叠。
\item 证明在这两个范畴上施行「取 Frobenius 迹」的操作，就得到所需的向量空间同构。
\item 对其他域上的光滑射影曲线、以及其他层论也表述几何 Langlands 对应，并证明所得的猜想本质上不依赖于所选定的设置。
\item 对特征 $0$ 的域上的曲线、在 D-模的层论中证明几何 Langlands 猜想。
\end{enumerate}
在下一节，我们将给出几何 Langlands 等价的一般性介绍；在随后的各节中，将聚焦于其表述与证明的不同方面。

\begin{remark}
在最初的数域情形，这一策略在第一步就无可挽回地失效，因为不存在 $X_{\overline{\F}_q}$ 的类比。长期以来，人们认为几何 Langlands 纲领所研究的结构不适用于最初的数域情形。然而，Fargues（2025）提议在一个特殊的「曲线」——Fargues–Fontaine 曲线上研究几何 Langlands 猜想，它在某种意义上把 $p$ 进数 $\Q_p$ 几何化。结果表明，这给出 $p$ 进域上局部 Langlands 猜想的一个精确版本，参看 Fargues–Scholze（2024）。在实数上也有平行的故事，参看 Ben-Zvi–Nadler（2013）、Scholze（2024）。局部算术 Langlands 目前的许多研究都建立在几何 Langlands 的范式之上。然而，在整体数域情形仍然需要真正新的想法。我们注意到，所需同构似乎极不可能通过「在某个范畴等价上取 Frobenius 迹」的类比而得到\footnote{在数域与 3 维流形的类比中，这最终归结为：有限域上的函数域对应以圆周为底的 3 维流形纤维化，而数域只对应 3 维流形，没有这种纤维化。}。
\end{remark}

\srcpage{11}
\textbf{致谢}。本报告只不过触及了皮毛，我为其中众多的遗漏、不精确与未竟之处致歉。关于几何 Langlands 对应之证明的另一篇生动得多的综述，我推荐 Ben-Zvi（2026）。我很高兴感谢许多数学家，他们教给我大量关于 Langlands 纲领、特别是几何 Langlands 纲领的知识，他们的洞见与远见塑造了我的理解：Johannes Anschütz、David Ben-Zvi、Arthur-César le Bras、Ana Caraiani、Dustin Clausen、Vladimir Drinfeld、Matt Emerton、Laurent Fargues、Dennis Gaitsgory、Toby Gee、Linus Hamann、David Hansen、Eugen Hellmann、Tim Kuppel、Vincent Lafforgue、Gérard Laumon、Lucas Mann、David Nadler、Michael Rapoport、Sam Raskin、Timo Richarz、Juan Esteban Rodríguez Camargo、Yiannis Sakellaridis、Germán Stefanich、Akshay Venkatesh、Kari Vilonen、Cong Xue、Zhiwei Yun、Xinwen Zhu，以及许多其他人。特别感谢 Dennis Gaitsgory 与 Sam Raskin 在准备本报告期间回答了我的许多问题；对于任何可能的误述，我谨致歉意。

\section{几何 Langlands 等价}\srcpage{11}

本节解释几何 Langlands 等价的大致形态。

作为一般设定，我们稍微改变记号，固定代数闭域 $k$ 上的一条光滑、紧合、连通曲线 $X$。为方便起见，固定 $X$ 上的一个 Theta 特征 $\Theta$，即 $\Theta^{\otimes 2} = \OmX$；这将是归一化几何 Langlands 等价所需要的。此外，固定 $k$ 上的一个约化群 $G$。由于 $k$ 代数闭，$G$ 总是分裂的；我们固定一个钉扎（pinning），特别地固定一个 Borel 子群与极大环面 $T \subset B \subset G$。

还有第三个选择——$k$ 上（有限型）概形的上同调理论，系数取在某个特征为 $0$ 的域 $e$ 中。\footnote{应当存在允许更一般系数的变体，但如同我们的主要参考文献，我们将只关注这一情形。}在现代的理解中，这由一个六函子形式化
\[
D \colon \Corr(\Sch_k^{\mathrm{ft}}) \to \PrL_e
\]
给出，即从对应的 $((2,1)$-)范畴到 $e$-线性可表现 $(\infty,1)$-范畴的松弛对称幺半函子；详见 Heyer–Mann（2024），其先声见 Gaitsgory–Rozenblyum（2017）、Liu–Zheng（2024）与 Mann（2022，附录 A.5）。

\begin{remark}
这里有两个不同的域 $k$ 与 $e$，我们将在二者之上做代数几何。最终我们会在 $D$-模的框架中工作，此时 $k = e$；但对函数域情形的应用，有 $k = \overline{\F}_q$，而 $e = \Ql$。
\end{remark}

\srcpage{12}目前还不知道如何对任意（合理的）上同调理论的选择表述几何 Langlands 等价。相反，对标准类型的上同调理论，例如 de Rham、Betti 与 étale，人们可以写出表述，而这些表述实际上各自用到了给定情形的某些特殊性质。出于这个原因，我不再过多停留于一般性；读者应把 $D$ 想成标准上同调理论之一：de Rham、Betti 或 étale。

第一个主要角色是 $X$ 上 $G$-丛的模空间 $\Bun_G$。它是光滑 Artin 叠，其 $S$-值点是 $X \times_k S$ 上 $G$-丛的群胚，其中 $S$ 为 $k$ 上任意概形。用下降把六函子形式化 $D$ 从概形延拓到 Artin 叠，便得到 $\Bun_G$ 上层的 $(\infty,1)$-范畴 $D(\Bun_G)$。

几何 Langlands 纲领的出发点是这样一个想法：这个范畴是自守形式向量空间的一个化身。这一想法源于 Drinfeld（1983）关于算术 Langlands 纲领函数域情形的工作。也就是说，对于 $\F_q$ 上的概形 $S$，构造 $S(\F_q)$ 上有趣函数的一般配方如下。从 $S_{\overline{\F}_q}$ 上一个带 Frobenius 结构的 $\ell$-进层 $\mathcal{F}$ 出发，即一个映射 $\phi \colon \mathrm{Fr}^*\mathcal{F} \to \mathcal{F}$。然后对每个 $\F_q$-点 $s \in S$，联系上 $\mathcal{F}_s$ 上 Frobenius 的迹
\[
\tr(\phi_s|\mathcal{F}_s) \in \Ql,
\]
其中 $\overline{s}$ 是 $s$ 上方某个几何点。事实上，最近人们注意到，Frobenius 迹这一概念在范畴层面也有意义，于是可以取
\[
\tr(\mathrm{Fr}|D(\Bun_G)) \in D(\Ql),
\]
得到一个（复形的）$\Ql$-向量空间。对层取 Frobenius 迹的通常操作给出一个映射
\[
\tr(\mathrm{Fr}|D(\Bun_G)) \to \cA_c(\Bun_G(\F_q)).
\]

下面的定理由 Arinkin、Gaitsgory、Kazhdan、Raskin、Rozenblyum 与 Varshavsky（2021）证明，它给出了「$D(\Bun_G)$ 是自守形式空间的范畴化」这一想法的一个非常强的体现：

\begin{theorem}
给 $D(\Bun_G)$ 取一个合适的定义（即取 étale $\Ql$-层构成的范畴），自然映射
\[
\tr(\mathrm{Fr}|D(\Bun_G)) \to \cA_c(\Bun_G(\F_q))
\]
是同构。
\end{theorem}

这个定理的一个关键部分是对 $D(\Bun_G)$ 的正确定义。遵循 Laumon（1987，第 6 节）的想法，必须取具有幂零奇异支撑的层，这里的幂零奇异支撑是就层的微局部理论而言的。我们将在下文第 3 节讨论这一点。

\srcpage{13}
\begin{remark}
这大概是我们可以看到的第一个场合，其中人们在几何 Langlands 纲领中做「用范畴做代数」。换句话说，大约一个世纪以前，人们逐步从用数做代数（加、乘……）走向用模做代数（直和、张量积……）；而在这里，人们做的是用范畴做代数（范畴同样容许直和、张量积……）。这样做的可能性依赖于 Lurie 的基础性工作，特别是 Lurie（2017），并要求全程以 $(\infty,1)$-范畴的方式工作。

几何 Langlands 猜想的证明展现了在用范畴做代数方面的极高技巧——所有关键步骤都在于对范畴的巧妙定义与操作（当然其间穿插着几何洞见）。值得注意的是，并且有点反直觉的是，这些操作最终能证明非常具体的东西，例如 Arthur–Ramanujan 猜想。
\end{remark}

几何 Langlands 纲领的目标是用「谱的语言」描述 $D(\Bun_G)$。粗略地说，人们遵循算术情形中的纲要：即第一步是构造一个作用在 $D(\Bun_G)$ 上的大交换算子代数。\footnote{延续上一条注，我们在这里使用「范畴中作用于一个范畴的交换代数」这一概念。注意，「范畴中的交换代数」与对称幺半范畴是同一回事。作为要记住的基本例子，对任意概形或叠 $X$，其（导出 $(\infty,1)$-）拟凝聚层范畴 $D_{\qc}(X)$ 是一个「范畴中的交换代数」；而对 $X$ 上的任意叠 $Y$，$D_{\qc}(Y)$ 成为 $D_{\qc}(X)$ 上的一个模。在 Lurie 的基础中，这并不像看起来那么可怕：首先，对任意对称幺半 $(\infty,1)$-范畴 $\mathcal{C}$，有（交换）代数 $A \in \mathcal{C}$ 的概念，以及 $A$-模 $M \in \mathcal{C}$ 的概念。这些基本上是显然的概念——$A$ 带有一个单位 $1 \to A$、一个乘法 $A \otimes A \to A$，并满足结合律与交换律。类似地，$A$-模 $M$ 带有一个作用映射 $A \otimes M \to M$，使通常的图表交换。重要的微妙之处在于，要求结合律、交换律等的图表不是条件，而是额外的数据，它们服从更高阶的相容性。现在 $\PrL_e$ 本身就是一个对称幺半 $(\infty,1)$-范畴，因此可以把上述概念也应用于 $\PrL_e$ 中的交换代数，得到其上的模等概念。}

事实上，人们可以用完全类似的方式定义 Hecke 算子。也就是说，对每个点 $x \in X$，有 Hecke 对应
\[
\begin{tikzcd}
& \Hck_{G,x} \arrow[dl, "p_1"'] \arrow[dr, "p_2"] & \\
\Bun_G & & \Bun_G,
\end{tikzcd}
\]
其中 $\Hck_{G,x}$ 参数化两个 $G$-丛以及它们在 $x$ 之外的同构。在算术情形中，Hecke 算子由元素 $h \in \mathbb{C}[G(\cO_x)\backslash G(F_x)/G(\cO_x)]$ 列举，并由如下操作给出。取 $\Bun_G(\F_q)$ 上的一个函数，把它拉回到 $\Hck_{G,x}(\F_q)$，乘以 $h$，然后沿 $p_2$ 的纤维求和。所有这些操作在层的六种运算下都有类比：拉回是 $p_1^*$，乘法\srcpage{14}是 $\otimes$，沿纤维求和是 $p_{2!}$。\footnote{事实上，在取 Frobenius 迹的层-函数对应下，这一类比是精确的。}

然而，还有一个问题有待回答：$h$ 的类比物是什么？这里需要的不是函数，而是一个层，且它应当是空间 $G(\cO_x)\backslash G(F_x)/G(\cO_x)$ 上的层，其中 $\cO_x \cong k[[t]]$。实际上，正如 $\Bun_G$，这应当被理解为一个模问题，即局部 Hecke 叠 $\Hck^{\loc}_G$，它是把 $k$-代数 $A$ 映到
\[
G(A[[t]])\backslash G(A((t)))/G(A[[t]])
\]
的函子的层化。采用记号 $\LG(A) = G(A((t)))$ 与 $\LpG(A) = G(A[[t]])$ 分别表示圈群与正圈群，这由
\[
\Hck^{\loc}_G = \LpG\backslash \LG/\LpG = \LpG\backslash \Gr_G
\]
给出，其中 $\Gr_G = \LG/\LpG$ 是称为仿射 Grassmann 簇的归纳投射概形。

正如通常的 Hecke 代数实际上是卷积下的代数，$e$-线性 $(\infty,1)$-范畴 $D(\Hck^{\loc}_G)$ 具有代数结构，即配备了一个幺半张量积。

几何 Langlands 纲领的一个关键输入是几何 Satake 等价，即 Satake 同构的范畴版本。为了陈述它，人们要过渡到反常层的子范畴。同样，这对一般的上同调理论没有明确含义，但在 de Rham、Betti 与 étale 情形中有既定含义。\footnote{然而，最近 Richarz–Scholbach（2021）的工作给出了几何 Satake 等价的一个动机版本，其办法是巧妙利用相关几何中出现的受限类型的动机——仅混合 Tate 动机。}

\begin{theorem}[Mirković–Vilonen，2007]
存在幺半范畴的自然等价
\[
\Perv(\Hck^{\loc}_G) \cong \Rep(\Gd),
\]
其中 $\Rep(\Gd)$ 是 Langlands 对偶群 $\Gd$ 的有限维表示的 Abel 范畴。
\end{theorem}

特别地，$\Perv(\Hck^{\loc}_G)$ 竟然奇迹般地是对称幺半的，这反映了通常 Hecke 代数的不显然的交换性。

\begin{remark}
人们可以在这一等价上取 Frobenius 迹，以还原通常的 Satake 同构。事实上，对于分歧群——即 $k((t))$ 上不承认 $k[[t]]$ 上约化模型的群 $G$——存在几何 Satake 等价的各种版本，可用以（重新）证明分歧群的 Hecke 代数的 Satake 同构，参见 Zhu（2024）、van den Hove（2024，定理 1.4）。\srcpage{15}这可以看作用来识别自守函数空间的那个论证的一个极小的玩具版本。
\end{remark}

我们注意到，几何 Satake 等价是几何给出比算术更精细之物的第一个关键点。也就是说，通常的 Satake 同构只把对偶群作为一个方便的记账工具引入——真正被 Satake 同构强制出来的只有 $\check T/W$。而在这里，人们可以用 Tannakian 理论，从仅涉及 $G$ 的数据真正重构出对偶群 $\Gd$。可以说，这实际上就是 Langlands 对偶群的正确界定——它就是在几何 Satake 等价中出现的 Tannakian 群。

\begin{remark}
几何 Satake 等价的证明，是先构造 $\Perv(\Hck^{\loc}_G)$ 上的交换性约束，即把幺半结构提升为对称幺半结构。这用到另一个关键想法：融合。上面我们在任意点 $x \in X$ 处引入了局部 Hecke 叠，但人们同样可以在任意有限点集 $(x_1,\ldots,x_n)$ 处定义局部 Hecke 叠。一般地，当这些点互不相同时，它就是各个点处局部 Hecke 茎的乘积。然而，当点碰撞时，相应的因子也碰撞：若所有点都是同一个点 $x$，它就只是 $x$ 处的局部 Hecke 叠。总而言之，这给出族中的一个（Beilinson–Drinfeld）局部 Hecke 叠
\[
\Hck^{\loc}_{X^n} \to X^n,
\]
它参数化点 $(x_1,\ldots,x_n) \in X^n$ 以及 $X$ 在诸 $x_i$ 之并处的完备化上的两个 $G$-丛，并带有它们在诸 $x_i$ 之外的同构。给定局部 Hecke 叠上的反常层 $\mathcal{F}_1,\ldots,\mathcal{F}_n$，人们可以在 $\Hck^{\loc}_{X^n}$ 上定义它们的融合积 $\mathcal{F}_1 * \cdots * \mathcal{F}_n$，它一般地只是诸 $\mathcal{F}_i$ 的外张量积。为把它延拓到处处，结果表明，「在 $X^n$ 上泛局部无环」这一性质是一种「平坦性」或「连续变化」假设，它保证延拓唯一，从而给出融合积。限制到对角点 $(x,\ldots,x) \in X^n$，这就还原了幺半张量积，但现在它显然是对称的。\footnote{一个技术要点是，交换性约束必须被某些符号修正，才能使几何 Satake 等价成立——这里的问题是在奇上同调次数中出现的通常的 Koszul 符号。}
\end{remark}

这一过程——让 $x \in X$ 在曲线上移动，乃至让两个点移动并碰撞——是通常算术 Langlands 纲领中看不到的新可能性。一言以蔽之，正是这种额外的灵活性使几何 Langlands 纲领更易接近。

利用几何 Satake 等价，对每个点 $x \in X$ 我们得到 $\Rep(\Gd)$ 在 $D(\Bun_G)$ 上的作用，从而得到 $\bigotimes_{x \in X}\Rep(\Gd)$ 的作用。这是\srcpage{16}无穷维 Hecke 代数
\[
\mathbb{T} = \bigotimes_p \mathbb{C}[G(\Z_p)\backslash G(\Q_p)/G(\Z_p)] \otimes Z(U(\gc))
\]
作用于自守函数向量空间的类比。

但正如我们刚刚注意到的，在这种几何情形中我们可以让点 $x \in X$ 变动。这给出族中的 Hecke 对应
\[
\begin{tikzcd}
& \Hck_{G,X} \arrow[dl] \arrow[dr] & \\
\Bun_G & & \Bun_G \times X,
\end{tikzcd}
\]
其中 $\Hck_{G,X}$ 参数化一个点 $x \in X$ 以及两个在 $x$ 之外同构的 $G$-丛。Hecke 算子现在可以定义为函子
\[
D(\Bun_G) \to D(\Bun_G \times X),
\]
对所有 $x \in X$ 插值通常的 Hecke 算子。这给出各个 $x \in X$ 的作用之间的一些非平凡相容性，并可用于产生谱作用。

为了解释谱作用，我们必须引入几何 Langlands 纲领中的另一个关键角色：$X$ 上 $\Gd$-局部系统的叠 $\LocSys_{\Gd}$，详见第 4 节。这是 $e$ 上的一个（导出）叠，参数化按所选上同调理论意义下 $X$ 上的 $\Gd$-局部系统。换言之，若所选上同调理论是 de Rham（分别地，Betti、étale），则 $\LocSys_{\Gd}$ 就是 $X$ 上 de Rham（分别地，Betti、étale）$\Gd$-局部系统的模空间。出于技术原因，必须把它视为导出代数几何中的一个对象。

有了这些，第一个定理（不是历史上第一个，而是在几何 Langlands 证明的逻辑上第一个）是谱作用的存在性。我们将在第 5 节进一步说明。

\begin{theorem}
对 $D(\Bun_G)$ 取合适的定义，存在 $D_{\qc}(\LocSys_{\Gd})$ 在 $D(\Bun_G)$ 上的一个作用，它细化 Hecke 作用。特别地，对所有 $V \in \Rep(\Gd)$ 与 $x \in X$，诱导函子
\[
D(\Bun_G) \to D(\Bun_G)
\]
由 $D_{\qc}(\LocSys_{\Gd})$ 中如下对象的作用给出：该对象对每个 $\Gd$-局部系统，联系上表示 $V$ 所给出的相伴局部系统在 $x \in X$ 处的纤维。
\end{theorem}

\begin{remark}
如果我们更仔细地写出谱作用与 Hecke 作用的相容性，这实际上会唯一确定谱作用。

这相当于说，自守形式空间从 $\mathbb{T}$-模提升为 $\LocSys_{\Gd}$ 上的模！因此，它已经给出了「自守到 Galois」的方向。
\end{remark}

\srcpage{17}
\begin{remark}
在函数域情形，Lafforgue（2018）对自守形式空间证明了这一陈述的类比：自守形式空间中出现的每个 Hecke 特征值系统，都必须来自 $X$ 上的一个 $\Gd$-局部系统。他的工作甚至允许分歧的情形。他的论证依赖于非常类似的要素——几何 Satake 等价，以及让曲线上的点移动的可能性；我们将在下文注 5.4 中描述他的结果的一个视角。V. Lafforgue 的工作把几何 Langlands 纲领与算术 Langlands 纲领重新统一起来，并直接启发了许多后续工作。

在 $D$-模情形中，谱作用的存在性需要沉重的非形式输入，这归功于 Gaitsgory（2015），他依赖于 Kac–Moody 局部化的结果，而这些结果直到最近才由 Arinkin、Beraldo、Campbell、Chen、Færgeman、Gaitsgory、Lin、Raskin 与 Rozenblyum（2024）写出。在 Betti 情形，正确的范畴 $D(\Bun_G) = D_{\Nilp}(\Bun_G)$ 由 Ben-Zvi 与 Nadler（2018）分离得到，而 Nadler 与 Yun（2019）在这一情形构造了谱作用。有点令人惊讶的是，在该设定中这本质上是形式废话；人们除了（高度结构化的）Hecke 作用之外不需要任何其他输入。他们的工作由 Arinkin、Gaitsgory、Kazhdan、Raskin、Rozenblyum 与 Varshavsky（2020b）推广到了 étale 情形。
\end{remark}

正如我们所讨论的，猜想 1.5 可以看作有两部分：首先，它预测自守形式空间提升为 $\LocSys_{\Gd}$ 上的一个模；然后它又预测这个模对应于一个具体的对象，即对偶化层。一般地，这是一个秩 $1$ 自由模。

在几何情形，定理 2.7 已给出了前半部分！现在人们希望把 $D(\Bun_G)$ 描述为 $D_{\qc}(\LocSys_{\Gd})$ 上的模。如同算术情形，最乐观的希望是 $D(\Bun_G)$ 成为 $D_{\qc}(\LocSys_{\Gd})$ 上的秩 $1$ 自由模；但这并不完全正确。相反，人们必须把由完美复形紧生成的 $D_{\qc}$ 替换为具有幂零奇异支撑的凝聚复形的归纳完备化，即 $\IndCoh_{\Nilp}(\LocSys_{\Gd})$，参见 Arinkin–Gaitsgory（2015）。

\begin{conjecture}[几何 Langlands 猜想]
存在 $D_{\qc}(\LocSys_{\Gd})$-线性 $(\infty,1)$-范畴的等价
\[
D(\Bun_G) \simeq \IndCoh_{\Nilp}(\LocSys_{\Gd}).
\]
\end{conjecture}

事实上，在这一情形中可以具体指定这个同构的归一化。也就是说，它应当把右端的结构层送到左端的 Whittaker 层。Whittaker 层的定义是 Poincaré 级数的直接类比：人们考察对应
\[
\begin{tikzcd}
& \Bun_U^{\Theta} \arrow[dl] \arrow[dr] & \\
\Bun_G & & \A^1,
\end{tikzcd}
\]
\srcpage{18}从 $\A^1$ 上的指数局部系统出发（即指数 $D$-模，分别地 Artin–Schreier 层；在 Betti 情形须作小的调整），把它拉回到 $\Bun_U^{\Theta}$，然后取到 $\Bun_G$ 的 $!$-前推。

事实证明，若存在所需形式的等价，它由这一要求唯一决定，参见下文注 9.4。再次，这相对于自守形式的情形是一个改进，见证了范畴陈述的额外刚性。

Gaitsgory–Raskin（2025b）、Arinkin、Beraldo、Campbell、Chen、Færgeman、Gaitsgory、Lin、Raskin 与 Rozenblyum（2024）、Campbell、Chen、Færgeman、Gaitsgory、Lin、Raskin 与 Rozenblyum（2024）、Arinkin、Beraldo、Chen、Færgeman、Gaitsgory、Lin、Raskin 与 Rozenblyum（2024）以及 Gaitsgory–Raskin（2026）这一系列论文，证明了特征 $0$ 的几何 Langlands 猜想。

\begin{theorem}[几何 Langlands 定理，特征 $0$]
若 $k$ 的特征为 $0$，则几何 Langlands 猜想成立，无论在 de Rham、Betti 还是 étale 情形中。
\end{theorem}

人们几乎可以消除对特征的假设：

\begin{theorem}[几何 Langlands 定理，特征 $p$，Gaitsgory–Raskin，2025a]
若 $k$ 的特征为正，且大于由 $G$ 决定的某个常数，\footnote{特别地，$p = \operatorname{char}(k)$ 不应是 $G$ 的挠素数。}则存在开闭子集 $\LocSys'_{\Gd} \subset \LocSys_{\Gd}$ 以及 $D_{\qc}(\LocSys_{\Gd})$-线性 $(\infty,1)$-范畴的等价
\[
D(\Bun_G) \simeq \IndCoh_{\Nilp}(\LocSys'_{\Gd}).
\]
\end{theorem}

正如之前所述，人们当然期望 $\LocSys'_{\Gd} = \LocSys_{\Gd}$；这对 $G = \mathrm{SL}_n$ 是已知的。

让我们以对证明策略的极简概述结束本节。首先，若 $k$ 的特征为 $0$，人们证明 de Rham、Betti 与 étale 情形中的陈述彼此等价；这是 Riemann–Hilbert 对应的一个应用，也是 Gaitsgory–Raskin（2025b）的主题。若 $k$ 的特征为正，人们通过把 $X$ 提升为 Witt 向量环 $W(k)$ 上的曲线，把 $X$ 形变到特征 $0$，并利用一般纤维上的结果，由退化推出 $X/k$ 上的结果；这是 Gaitsgory–Raskin（2025a）的主题。目前，这一论证还不能完全奏效，原因是混合特征下 étale 层的奇异支撑理论有微妙之处——这最终导致了开闭子集 $\LocSys'_{\Gd} \subset \LocSys_{\Gd}$ 的出现。

此后，证明在 de Rham 情形、$k$ 的特征为 $0$ 下进行。证明对（$G$ 的）秩作归纳，类比于把自守形式分解为 Eisenstein 级数的连续谱与尖点形式，这些 Eisenstein 级数由 $G$ 的真 Levi 子群诱导。\srcpage{19}至少有三个关键要素。

\begin{enumerate}
\item Whittaker 系数的保守性结果，归功于 Færgeman–Raskin（2025）。粗略地说，这保证 $D(\Bun_G)$ 上所有缓增的、例如尖点的层都能被检测到。
\item Kac–Moody 局部化。这给出一个非常一般的程序，从谱数据（即 opers）出发，在 $D(\Bun_G)$ 上产生（自守的）对象；它可追溯到 Beilinson–Drinfeld（1996）的开创性工作。这是仅在 $D$-模设定中存在的操作，并最终导致了「Galois 到自守」的方向。
\item $\LocSys_{\Gd}$ 的几何。对于最后的论证，即识别尖点部分，人们要对 $D$-模设定中 $\LocSys_{\Gd}$ 的几何作一些非常具体的分析。值得注意的是，与算术情形中尖点形式及其对应的 $L$-参数构成离散谱不同，在几何情形中，即便不可约 $\Gd$-局部系统的轨迹也是连通的，从而构成连续族，使插值论证得以进行。
\end{enumerate}

最后，对定理 2.12 中的等价取 Frobenius 迹，便可推出定理 1.7。这一步能还原自守形式空间，本身就是一个非平凡的陈述，参见第 8 节。

\begin{remark}
正如之前所述，几何 Langlands 猜想的一般形态，已被证明极具影响力，甚至超出了它最初的光滑紧合曲线设定。也就是说，在 Fargues–Scholze（2024）中，它在 Fargues–Fontaine 曲线上得到解释，与 $p$ 进域上的局部 Langlands 对应相联系；而在 Ben-Zvi–Nadler（2013）与 Scholze（2024）中，它在扭量 $\mathbb{P}^1$ 上得到解释（以略有不同的方式），与阿基米德域（即 $\R$ 与 $\C$）上的局部 Langlands 对应相联系。此外，$p$ 进域上局部 Langlands 对应的各种更精细形式——处理 $p$ 进向量空间上的 Banach 表示或局部解析表示——目前正用这一范式研究，例如参见 Emerton–Gee–Hellmann（2025）。这打破了该领域的一个僵局：在 $\mathrm{GL}_2(\Q_p)$ 之外，甚至猜想性的图景直到最近都还不清楚。
\end{remark}

\section{$\Bun_G$ 上的层}\srcpage{19}

为与我们由函数域 Langlands 猜想所驱动的主线保持一致，我们先在 étale 情形下定义 $\Bun_G$ 上正确的层范畴。因此，$k$ 是某个代数闭域（常见的是 $k=\overline{\F}_q$，但也可以允许 $k$ 的特征为 $0$），$e=\Ql$，而我们的层论为
\[
S\mapsto \Ind D^b_c(S,\Ql).
\]
\srcpage{20}其中 $D^b_c(S,\Ql)$ 表示可构造 étale $\Ql$-层的有界导出范畴。这里，我们把它视为稳定 $(\infty,1)$-范畴，再过渡到 Ind-对象，得到可表现 $\Ql$-线性 $(\infty,1)$-范畴。

六函子形式化的标准选择在于采用通常的张量 $\otimes$、拉回 $f^*$ 与紧支撑上同调 $f_!$-函子（及其右伴随）。在 Arinkin、Gaitsgory、Kazhdan、Raskin、Rozenblyum 与 Varshavsky（2020b）中作了不同的选择，即采用 !-张量 $\otimes^!$、例外拉回 $f^!$ 与通常上同调 $f_*$（及其右伴随）。注意
\[
D^b_c(S,\Ql)
\]
是自对偶的，即经由 Verdier 对偶等价于其反范畴；通常的 $\otimes$、$f^*$ 与 $f_!$-函子保持 $D^b_c(S,\Ql)$；而其余函子 $\otimes^!$、$f^!$ 与 $f_*$ 是 $\otimes$、$f^*$ 与 $f_!$ 在 Verdier 对偶下的共轭。因此，在紧对象层面，两种选择的差别本质上只是表面上的，仅仅在于施加一次 Verdier 对偶。然而，在过渡到 Ind-范畴并把形式化推广到叠之后，两种选择一般真正不同。幸运的是，在 $\Bun_G$ 上，这一选择仍然无关紧要：即这是光滑 Artin 叠，故 $D(\Bun_G)$ 经由光滑图谱的下降来定义。但对光滑映射 $f$，映射 $f^*$ 与 $f^!$ 至多相差一个平移后一致，所以所得的范畴与形式化的选取无关。

\begin{remark}
我必须承认，我觉得 $(-^!,-_*)$-形式化令人困惑，以至于对 Arinkin、Gaitsgory、Kazhdan、Raskin、Rozenblyum 与 Varshavsky（2020b、2020a、2021）这一系列论文，我只有表面的理解。然而，他们的选择也有清楚的动机：即与 D-模的六函子相容。

另一方面，就几何 Langlands 等价而言，这一选择并不重要。为了我个人方便，本报告将采用 $(-^*,-_!)$-形式化。
\end{remark}

\begin{remark}
与任一六函子形式化 $D$ 相联系的一个关键构造，是核的（对称幺半可表现 $(\infty,2)$-）范畴 $K_D$。定义它的方式有好几种；最简洁的一种是取 $D$ 在
\[
\Fun(\Corr(\Sch^{\mathrm{ft}}_k),\PrL_e)
\]
中的模（注意：松弛对称幺半函子恰是关于 Day 卷积的交换代数对象）。它由有限型 $k$-概形 $X$ 对应的对象 $[X]\in K_D$ 生成，且
\[
\Hom_{K_D}([X],[Y])=D(X\times_k Y),
\]
其中单位态射由 $\Delta_!\mathbf{1}$ 给出，复合由卷积给出。

在任一 2-范畴中，都有伴随态射的概念 —— 若存在态射 $\epsilon\colon\id_Y\to GF$（在 $\End(Y)$ 中）与 $\eta\colon FG\to\id_X$（在 $\End(X)$ 中），\srcpage{21}使得复合
\[
F\xrightarrow{F\epsilon}FGF\xrightarrow{\eta F}F,\qquad
G\xrightarrow{\epsilon G}GFG\xrightarrow{G\eta}G
\]
等价于恒等，则 $F\colon X\to Y$ 是 $G\colon Y\to X$ 的左伴随。伴随在存在时于同构意义下唯一。给定层 $A\in D(X)$，可将 $A$ 视为 $K_D$ 中的态射 $[*]\to[X]$，并询问左（分别右）伴随是否存在；这引出了「suave」（分别「prim」）层的概念；参见 Heyer 与 Mann（2024）。「suave」与经典意义上的泛局部无环概念密切相关，例如参见 Hansen 与 Scholze（2023）。\footnote{对常值层 $A=\mathbf{1}$ 而言，「suave」是光滑性的减弱，且与泛无环性相关，故其名字中的字母 s、u、a 甚为贴切。「prim」则是固有性的减弱。}
\end{remark}

\begin{remark}
尽管 $(-^*,-_!)$ 与 $(-^!,-_*)$ 两种形式化在 $(\infty,1)$-范畴层面大体一致，它们给出的核 $(\infty,2)$-范畴却很不相同；特别地，比较由此得到的伴随概念颇为微妙。值得注意的是，在下一个定理中，用 prim 性给出的刻画只适用于 $(-^*,-_!)$-形式化。
\end{remark}

对几何 Langlands 对应，必须限制到 $D(\Bun_G)$ 的某个子范畴。该子范畴由下述定理刻画。为陈述它，回忆：光滑概形或叠 $X$ 上层的奇异支撑是余切丛 $T^*X$ 的一个子集，衡量该层的奇异性。当 $X=\Bun_G$ 时，余切丛 $T^*X$ 成为 Higgs 丛的模空间 $\Higgs_G$，而幂零奇异支撑指的是相应的 Higgs 丛幂零，即位于 Hitchin 纤维化的零纤维中。对 étale 层，奇异支撑的概念由 Beilinson（2016）建立。

与 prim 对象的等价，其动机来自 Hansen--Mann 在 Fargues--Fontaine 曲线情形中的类似想法；它在 Kuppel 的硕士论文（2024，定理 3.3.24）中得到证明，但以 Arinkin、Gaitsgory、Kazhdan、Raskin、Rozenblyum 与 Varshavsky（2020b，猜想 14.1.8）为前提，而该猜想现已被 Gaitsgory 与 Raskin（2025a）证明。

\begin{theorem}[Arinkin、Gaitsgory、Kazhdan、Raskin、Rozenblyum 与 Varshavsky，2020b]
设 $p$ 相对于 $G$ 充分大；特别地，它对 $G$ 是非挠的素数。\footnote{这一假设是为了在 $\gc$ 中有良好的幂零锥概念。}设 $A\in D(\Bun_G)$。以下条件等价。

(1) 层 $A$ 具有幂零奇异支撑，即其奇异支撑落在 $T^*\Bun_G$ 的幂零锥中。

(2) 层 $A$ 是 Hecke-lisse 的，即对所有 $V\in\Rep(\Gd)$，Hecke 平移 $T_V(A)\in D(\Bun_G\times X)$ 沿 $X$-方向局部常值。

\srcpage{22}(3) 层 $A$ 属于
\[
\Ind D(\Bun_G)^{\prim}\subset D(\Bun_G),
\]
其中 $D(\Bun_G)^{\prim}$ 表示 prim 对象。

记
\[
D_{\Nilp}(\Bun_G)\subset D(\Bun_G)
\]
为满足这些等价条件的对象构成的全子范畴。
\end{theorem}

\begin{proof}[证明概要]
容易的方向是 (1) 与 (3) 各自推出 (2)。对 (1)，用到奇异支撑在光滑拉回与固有前推下表现良好，并借助关于幂零锥行为的一个计算来完成。参见 Nadler 与 Yun（2019，定理 5.2.1、6.1.1）。

对 (3)，先指出 prim 层自动是紧对象（Heyer 与 Mann，2024，引理 4.4.18(ii)），故该 Ind-范畴实际上是全子范畴。为验证 (3) 推出 (2)，便可归约到 prim 层。现在用到 prim 层在固有 $*$-拉回与 suave $!$-前推下稳定。这表明 $T_V(A)\in D(\Bun_G\times X)$ 必定仍是 prim 的。特别地，对任一可构造的 $B\in D(\Bun_G)$ —— 它此时是 suave 的，从而作为方向 $D(\Bun_G)\to D(*)$ 上的核具有右伴随 —— 对象 $p_!(B\otimes T_V(A))\in D(X)$ 是 prim 的。但对固有的 $X$，prim 性等价于可对偶性（Heyer 与 Mann，2024，推论 4.5.18），所以这是 $X$ 上的局部系统。若取任一闭点 $s\colon *\to\Bun_G$ 并对 $B=s_!\Ql$ 施以此论，便知 $s^*T_V(A)\in D(X)$ 是局部系统。在 $k$ 中作任意基变换后再施以同样的论证，表明 $T_V(A)$ 在 $\Bun_G$ 各点处的所有茎都是 $X$ 上的局部系统；而这正是 (2)。

由 (2) 到 (1) 的逆向蕴涵在 Arinkin、Gaitsgory、Kazhdan、Raskin、Rozenblyum 与 Varshavsky（2020b，定理 14.4.4）中建立。大致说来，要证明的是：若存在某个非幂零的奇异支撑，则通过适当选取的 Hecke 算子，会把它显现为曲线方向上的某个非平凡奇异支撑。

至于 (2) 到 (3) 的蕴涵，参见第 7 节。
\end{proof}

最后，我们说明在 Betti 与 de Rham 情形中使用的是 $\Bun_G$ 上的哪些层范畴。Betti 情形曾在 Ben-Zvi 与 Nadler（2018）中考虑。此处 $k=\C$，层论为
\[
D(S)=D(S(\C),e),
\]
即局部紧拓扑空间 $S(\C)$ 上 $e$-向量空间层的复形。同样取 $D_{\Nilp}(\Bun_G)\subset D(\Bun_G)$，即具有幂零奇异支撑的 Betti 层构成的全子范畴。同样可以证明它等价于 Hecke-lisse；但在 Betti 情形，它并不蕴涵 (ind-)prim。作为玩具例子，\srcpage{23}把 $\Bun_G$ 换成 $S^1$，并以零奇异支撑代替幂零奇异支撑：\footnote{正如我们将在第 6 节讨论的，$G=\Gm$ 的情形实质上就产生这一情形。}于是得到由 $S^1$ 上局部常值层构成的全子范畴
\[
D(*/\Z)\subset D(S^1).
\]
现在 $D(*/\Z)\cong D_{\qc}(\mathbb{G}_{m,e})$ 等价于 $\mathbb{G}_{m,e}$ 上的拟凝聚层 —— 两边都是带一个自同构的 $D(e)$ 的对象。$\mathbb{G}_{m,e}$ 的结构层给出一个紧对象，它在 $*/\Z$ 上、从而通过拉回在 $S^1$ 上，是茎为无穷维的层；这不是 (ind-)prim 的。

Betti 与 étale 情形的差异，是六函子形式化精确定义的产物：$\Ind D^b_c(-,\Ql)$ 的类比物毋宁是 $\Ind D_{\mathrm{cons}}(-,e)$，即 Zariski-可构造 Betti 层的 Ind-范畴。这一选择将与 étale 情形完全类似。

de Rham 情形更容易：此处 $k$ 是任意的特征 $0$（代数闭）域，系数 $e=k$ 与 $k$ 一致，并取 $\Bun_G$ 上 D-模的完整导出范畴 $D(\Bun_G)$。值得注意的是，并不过渡到 (ind-)holonomic D-模的某个子范畴。

\section{局部系统的叠}\srcpage{23}

在 Langlands 对偶一侧，主角是局部系统的叠。本节讨论它在不同情形下的定义。

此处最熟悉的情形是 Betti 情形。于是 $X$ 是 $k=\C$ 上的光滑射影曲线，等价地说，是紧 Riemann 面。为简单起见，我们排除 $\mathbb{P}^1$ 的情形，并假设 $X$ 的亏格 $g\ge 1$。那么作为拓扑空间，$X$ 是其基本群 $\pi_1=\pi_1(X)$ 的 $K(\pi_1,1)$，而该群是有限表现群 —— 事实上，它有著名的表现
\[
\pi_1=\langle a_1,b_1,a_2,b_2,\dots,a_g,b_g\mid [a_1,b_1][a_2,b_2]\cdots[a_g,b_g]=1\rangle,
\]
有 $2g$ 个生成元与一个关系。

我们可以用任意特征 $0$ 的代数闭域 $e$ 作系数；读者可假设 $e=\C$，但 $e=\Ql$ 也将是相关的（不过无论如何，它也同构于 $\C$）。在此情形，$X$ 上的 $\Gd$-局部系统无非由同态
\[
\pi_1\to\Gd(e)
\]
给出，至多差 $\Gd(e)$-共轭。更确切地说，有仿射概形 $\Hom(\pi_1,\Gd)$ 参数化同态 $\pi_1\to\Gd$，而局部系统的叠是所得的商叠
\[
\LocSys^{\Betti}_{\Gd}=\Hom(\pi_1,\Gd)/\Gd.
\]

\srcpage{24}具体地，利用 $\pi_1$ 的表现，有
\[
\Hom(\pi_1,\Gd)=\Gd^{2g}\times_{\Gd}*,
\]
它是映射
\[
\Gd^{2g}\to\Gd\colon (a_1,b_1,\dots,a_g,b_g)\mapsto [a_1,b_1]\cdots[a_g,b_g]
\]
在 $1\in\Gd$ 处的纤维。

我们指出，一般情形下 —— 例如当 $\Gd$ 是环面时 —— 这个映射不是平坦的，而纤维积
\[
\Hom(\pi_1,\Gd)=\Gd^{2g}\times_{\Gd}*
\]
必须在导出意义下取，作为导出仿射概形，其函数环是导出张量积
\[
\cO(\Gd^{2g})\otimes^L_{\cO(\Gd)}e.
\]

若 $\Gd$ 半单且 $g\ge 2$，这个问题不会出现，例如参见 Aizenbud 与 Avni（2016）。一般地，$\Hom(\pi_1,\Gd)$ 只有温和的奇点：正如该表现所示，它是完全交。

因此，在 Betti 情形，局部系统的叠是极为基本且被充分研究的对象：至少当 $\Gd=\mathrm{GL}_n$ 时，它不过是视为叠的特征簇。

若 $e=\C$，则 Riemann--Hilbert 对应给出局部系统与带平坦联络的向量丛之间的翻译。这就把我们带到局部系统叠的 de Rham 体现。它有相同的 $\C$-点集（更确切地说是群胚），但代数结构不同。更一般地，若 $k$ 是任一特征 $0$ 的（代数闭）域，可取 $e=k$，并定义代数叠
\[
\LocSys^{\dR}_{\Gd},
\]
其 $A$-值点（对 $k$-代数 $A$ 而言）是 $X\times_{\Spec(k)}\Spec(A)$ 上带 $X$-方向平坦联络的 $\Gd$-丛构成的群胚。等价地，它们是 $X_{\dR}\times_{\Spec(k)}\Spec(A)$ 上的 $\Gd$-丛，其中 $X_{\dR}$ 是 $X$ 的 Simpson de Rham 叠。回忆：这是 $X$ 按对角线的形式完备化所作的商 $X_{\dR}=X/(X\subset X\times X)^{\wedge}$，而 $X_{\dR}$ 上的拟凝聚层恰是 $X$ 上的 D-模。

遗忘联络给出代数叠的映射
\[
\LocSys^{\dR}_{\Gd}\to\Bun_{\Gd},
\]
它在 $\Gd$-丛 $E$ 处的纤维分类 $E$ 上的联络。$E$ 上联络的集合构成以仿射空间 $\Gamma(X,E\times^{\Gd}\check{\gc}\otimes\OmX)$ 为模型的拟挠子。同样，一切都必须在导出代数几何的世界中解释，故 $A$ 是 $D_{\ge 0}(k)$ 中的代数，等等。在这一解释下，$\LocSys^{\dR}_{\Gd}$ 仍是局部完全交 —— $\Bun_{\Gd}$ 光滑，而到它的投影是局部完全交。关键在于复形 $\Gamma(X,E\times^{\Gd}\check{\gc}\otimes\OmX)$ 集中在上同调次数 $0$ 与 $1$。

最后考虑 étale 情形；此时 $e=\Ql$，且可允许特征不同于 $\ell$ 的任意代数闭域 $k$。在此情形，有投射有限群 $\pi_1^{\et}=\pi_1^{\et}(X)$，当 $k=\C$ 时它是\srcpage{25}上文所述离散群 $\pi_1$ 的投射有限完备化。若 $k$ 有正特征 $p$，则通过选取整环 $\cO_C$ 上的光滑射影曲线 $X_{\cO_C}$（$\cO_C$ 为某剩余域为 $k$ 的完备代数闭 $p$ 进域 $C$ 的整数环），并使 $X_{\cO_C}$ 形变 $X=X_k$，得到特殊化映射
\[
\pi_1^{\et}(X_C)\to\pi_1^{\et}(X),
\]
它是满射，把 $\pi_1^{\et}(X)$ 写成离散群 $\pi_1$ 的投射有限完备化的商。

理想情形下，人们会考虑 $\pi_1^{\et}$ 的连续表示的某个模空间；然而这样的模空间充其量只能生活在刚性解析几何的领域。就我们目前的目的而言，我们希望得到通常（导出）代数几何中的模叠。

粗略地说，Arinkin、Gaitsgory、Kazhdan、Raskin、Rozenblyum 与 Varshavsky（2020b）采取的办法，是把 $\pi_1^{\et}$ 换成它在 $\Ql$ 上的 pro-代数完备化，并考虑其表示。在 Tannakian 视角下，这意味着考虑全子范畴
\[
\Lisse(X)\subset D^b_c(X,\Ql),
\]
即 lisse $\Ql$-层的有界复形，然后与其 Ind-范畴 $\Ind\Lisse(X)$ 打交道。于是对任一导出 $\Ql$-代数 $A$，可以考虑左 $t$-正合的 $\otimes$-函子
\[
\Rep(\Gd)\to\Ind\Lisse(X)\otimes_{\Ql}A;
\]
这定义了叠 $\LocSys^{\et}_X$（在文献中记作 $\LocSys^{\restr}_X$）。\footnote{对 $X=\mathbb{P}^1$，需要作轻微的修改。}

为理解这个模叠的性质，我们指出，在 Betti 情形也可以这样做。此时令
\[
\LocSys^{\coarse}_{\Gd}=\Hom(\pi_1,\Gd)/\!/\Gd=\Spec(\cO(\LocSys^{\Betti}_{\Gd}))
\]
为粗略模空间。那么 Betti 叠的限制版本是
\[
\LocSys^{\Betti}_{\Gd}\times_{\LocSys^{\coarse}_{\Gd}}\bigsqcup_{x\in\LocSys^{\coarse}_{\Gd}(\C)}(\LocSys^{\coarse}_{\Gd})^{\wedge}_{x},
\]
即在粗略模空间层面只允许无穷小变化。模叠 $\LocSys^{\et}_{\Gd}$ 则是连通分支的并，对应于那些使相应的 $\Gd$-表示 $\pi_1\to\Gd(e)=\Gd(\Ql)$ 连续延拓到 $\pi_1^{\et}$ 的 $x$。

\section{谱作用}\srcpage{25}

谱作用在各种情形下给出 $D_{\qc}(\LocSys_{\Gd})$ 在 $D(\Bun_G)$ 上的一个作用。为得到它，可以更一般地追问：给出 \srcpage{26}$D_{\qc}(\LocSys_{\Gd})$ 的一个作用，究竟需要什么。在 Betti 情形，这有一个漂亮的答案，归功于 Nadler 与 Yun（2019）：

\begin{theorem}
设 $C\in\PrL_{\e}$。在 $C$ 上给出 $D_{\qc}(\LocSys^{\Betti}_{\Gd})$ 的一个作用，等价于给出幺半 $e$-线性范畴之间的正合函子
\[
\Rep(\Gd^I)\to\End(C)^{B\pi_1^I},
\]
且关于有限集 $I$ 函子式地成立。这里，右端表示幺半范畴 $\End(C)$ 中的 $\pi_1^I$-等变对象，而 $\End(C)$ 是 $C$ 的 $e$-线性、保余极限自函子构成的范畴。
\end{theorem}

这里，我们用 $\Rep(\Gd)$ 表示 $\Gd$ 的有限维表示构成的小 Abel 范畴。当从它出发作映射时，我们隐含地用到：对可表现稳定的 $C$，正合函子 $\Rep(\Gd)\to C$ 与保余极限函子 $D_{\qc}(\ast/\Gd)\to C$ 是同一回事。

\begin{remark}
定理 5.1 有一个更初等的类似物：它描述的不是 $\LocSys^{\Betti}_{\Gd}$ 上的范畴层，而是模层，即拟凝聚层。这可以用同样的方式证明，结论是：对象 $M\in D_{\qc}(\LocSys^{\Betti}_{\Gd})$ 等价于函子
\[
\Rep(\Gd^I)\to D(\ast/\pi_1^I)=D(e)^{B\pi_1^I}
\]
的数据，且关于 $I$ 函子式。这把 $M$ 送到对应
\[
\begin{tikzcd}
\LocSys^{\Betti}_{\Gd}\times\ast/\pi_1^I \arrow[r] \arrow[d] & \ast/\Gd^I \\
\ast/\pi_1^I
\end{tikzcd}
\]
上以 $M$ 为核的函子。这里，上方的箭头伴随于
\[
\LocSys^{\Betti}_{\Gd}=\Map(\ast/\pi_1,\ast/\Gd)\to\Map(\ast/\pi_1^I,\ast/\Gd^I).
\]
\end{remark}

\begin{proof}[证明概要]
我们给出一个论证，所用的是 Gestalt（格式塔）的语言，取自 Scholze（2026）关于与 Stefanich 合作工作的课程。$e$ 上的一个 Gestalt $X$ 由一列 $(\cO(X)_0,\cO(X)_1,\dots)$ 给出。其中 $\cO(X)_0$ 是 $D(e)$ 中的交换环，即「$X$ 上的函数环」。$\cO(X)_1$ 是幺半 $e$-线性 $(\infty,1)$-范畴，其单位的自同态给出 $\cO(X)_0$，即「$X$ 上的拟凝聚层范畴」。$\cO(X)_2$ 是幺半 $e$-线性 $(\infty,2)$-范畴，其单位的自同态给出 $\cO(X)_1$，即「$X$ 上的线性范畴层构成的 2-范畴」。……因此，Gestalten 旨在捕捉所有范畴层级上的高阶代数。出人意料的结论是：它们本质上构成一个 $(\infty,1)$-拓扑斯，从而可以像几何对象一样被思考与处理。因此，Gestalten 的语言把高阶范畴中困难的代数，翻译为更易驾驭的几何。这与概形的语言把困难的交换代数翻译为更易驾驭的几何，方式相同。

\srcpage{27}在 $e$ 上的 Gestalten 中，考虑映射 Gestalt
\[
\Map(\ast/\pi_1,\ast/\Gd)=\Hom(\pi_1,\Gd)/\Gd.
\]
这实际上就是 $\LocSys^{\Betti}_{\Gd}$，经由从代数叠到 Gestalten 的自然函子而视为一个 Gestalt。另一方面，映射 Gestalt 有一个一般公式。即：若 $X$ 是 $n$-suave 的、$Y$ 是 $n$-affine 的，则映射 Gestalt $\Map(X,Y)$ 自身是 $n$-affine 的，从而对应某个对称幺半可表现 $(\infty,n)$-范畴
\[
\cO(\Map(X,Y))_n\in\operatorname{CAlg}(\mathrm{nPr}_e);
\]
这里，$\mathrm{nPr}_e$ 表示可表现 $e$-线性 $(\infty,n)$-范畴构成的对称幺半范畴。此时 $\cO(\Map(X,Y))_n$ 实际上是万有的 $C\in\operatorname{CAlg}(\mathrm{nPr}_e)$，它配有对称幺半函子
\[
\cO(Y)_n\to\Hom(\cO(X)^!_n,C).
\]
这可由 $\Map(X,Y)$ 的定义以及 suave 映射的基变换公式推出（Scholze，2026，命题 6.18），后者涉及 $\cO(X)_n$ 的一个前对偶 $\cO(X)^!_n$。若忘掉 $C$ 上的交换代数结构，则万有的 $C$ 会是 $\cO(Y)_n\otimes\cO(X)^!_n$。回忆：Lurie 高阶代数中交换代数的概念，以有限集的操作子来编码。因此，万有的对称幺半 $C$ 可写成如下函子的余端：它把任一有限集 $I$ 送到
\[
\cO(Y)^{\otimes I}_n\otimes(\cO(X)^!_n)^{\otimes I}.
\]

在我们的情形，$X=\ast/\pi_1$ 是 $0$-étale 的，从而对任意 $n\ge 0$ 都是 $n$-suave 的；而 $Y=\ast/\Gd$ 是 $1$-affine 的。因此，该映射 Gestalt 是 $1$-affine 的。事实上，我们已知它由代数叠 $\LocSys^{\Betti}_{\Gd}=\Hom(\pi_1,\Gd)/\Gd$ 给出，所以
\[
\Map(\ast/\pi_1,\ast/\Gd)=\mathrm{Gest}(D_{\qc}(\LocSys^{\Betti}_{\Gd})).
\]
特别地，由于它是 $1$-affine 的，有
\[
\cO(\Map(\ast/\pi_1,\ast/\Gd))_2=\operatorname{Mod}_{D_{\qc}(\LocSys^{\Betti}_{\Gd})}(\mathrm{1Pr}_e),
\]
而这正是我们想理解的。按上述一般描述，它可描述为如下函子的余端：
\[
I\mapsto\cO(\ast/\Gd)^{\otimes I}_2\otimes(\cO(\ast/\pi_1)^!_2)^{\otimes I}.
\]
这里，$\cO(\ast/\Gd)^{\otimes I}_2=\operatorname{Mod}_{\Rep(\Gd)^{\otimes I}}(\mathrm{1Pr}_e)$，而加上第二个张量因子后，它就成为其中的 $\pi_1^I$-等变对象。拆解开来即得定理。
\end{proof}

\srcpage{28}定理 5.1 的数据几乎就是 Hecke 作用给我们的东西！更确切地说，对任一有限集 $I$ 与 $X^I$ 中任意的 $I$ 个点 $(x_i)_{i\in I}$，我们有 Hecke 对应
\[
\begin{tikzcd}
& \Hck_{G,(x_i)_i} \arrow[ld] \arrow[rd] & \\
\Bun_G & & \Bun_G.
\end{tikzcd}
\]
利用 $I$ 个点处几何 Satake 等价的一个版本，我们得到 $e$-线性幺半函子
\[
\Rep(\Gd^I)\to\Perv(\Hck_{G,(x_i)_i}),
\]
它关于目标上的卷积结构是幺半的。另一方面，以 Hecke 叠上的反常层作核，我们得到 $e$-线性幺半函子
\[
\Perv(\Hck_{G,(x_i)_i})\to\End(D(\Bun_G)).
\]
复合，便得到 $e$-线性幺半函子
\[
\Rep(\Gd^I)\to\End(D(\Bun_G)).
\]
但现在我们可以让点 $(x_i)_{i\in I}\in X^I$ 移动。若上述函子在点变动时适当地局部常值，就会给出所要的映射
\[
\Rep(\Gd^I)\to\End(D(\Bun_G))^{B\pi_1^I},
\]
即一个 $\pi_1^I=\pi_1(X^I)$-作用。

更确切地说，Hecke 作用实际上给出函子
\[
\Rep(\Gd^I)\times D(\Bun_G)\to D(\Bun_G\times X^I).
\]
在目标上，我们有沿 $X^I$-方向局部常值的层构成的全子范畴。它们等价地就是
\[
D(\Bun_G)^{B\pi_1^I}\cong D(\Bun_G\times\ast/\pi_1^I)\to D(\Bun_G\times X^I)
\]
的本质像中的层，这里用的是沿 $X^I\to\ast/\pi_1^I$ 的拉回。

正如我们在定理 3.4 处所讨论的，这对 $\Bun_G$ 上的所有层并不成立，但对
\[
D_{\Nilp}(\Bun_G)\subset D(\Bun_G)
\]
确实成立。在 étale 情形，情况完全平行；定理 5.1 有直接的类比成立，参见 Arinkin、Gaitsgory、Kazhdan、Raskin、Rozenblyum 与 Varshavsky（2020b，主定理 8.1.4）。

\begin{corollary}[Nadler 与 Yun，2019；Arinkin、Gaitsgory、Kazhdan、Raskin、Rozenblyum 与 Varshavsky，2020b]
在 Betti 或 étale 情形，存在 $D_{\qc}(\LocSys_{\Gd})$ 在 $D_{\Nilp}(\Bun_G)$ 上的谱作用（其中一个来源见主定理 14.3.2）。
\end{corollary}

\begin{remark}
\srcpage{29}正如我们在猜想 1.5 之后所讨论的，这给出几何 Langlands 的「自守到 Galois」方向。Lafforgue（2018）用非常类似的想法，建立了函数域上算术 Langlands 的「自守到 Galois」方向。更确切地说，本注中仍设 $X$ 是 $\Fq$ 上的曲线。\footnote{与几何 Langlands 纲领的当前状况不同，对下面的讨论我们甚至可以允许分歧，正如 V. Lafforgue 所做的那样。}于是我们可以把 $\cA_c(\Bun_G(\Fq))$ 提升为（算术的）$\LocSys_{\Gd}$ 上的模，另参见 Lafforgue 与 Zhu（2019）。由注 5.2 的一个版本，这等价于给出函子
\[
\Rep(\Gd^I)\to D(\Ql)^{BW^I}
\]
其中 $W$ 表示有限域上该曲线的 Weil 群，即几何基本群被 $\Z$ 的半直积。这恰好由取 shtuka 空间的上同调这一函子给出！下面给出 shtuka 模空间的一个稍显非标准的表述，它却简化了关于「部分 Frobenius」的讨论。回忆：直观上，人们把映射
\[
X\to\Spec(\Fq)
\]
想成一个纤维化于圆周之上的三维流形。有一个简便的办法把它变成真正的数学。我们令
\[
\Spec(\Fq)^{\Weil}=(\Spec(\overline{\F}_q)\times\R_{>0})/(\Frob_{\Fq}\times q)^{\Z},
\]
其中 $\R_{>0}$ 作为凝聚态集进入（即作为 pro-étale 代数空间，把它写成局部投射有限集的商）。于是 $\Spec(\Fq)^{\Weil}$ 的底层拓扑空间是一个圆周 $S^1=\R_{>0}/q^{\Z}$，所有剩余域都是 $\overline{\F}_q$；但绕圆周走一圈，就作用以 Frobenius 自同构。类似地，可以定义
\[
X^{\Weil}=(X_{\overline{\F}_q}\times\R_{>0})/(\Frob_{\overline{\F}_q}\times q)^{\Z}\to\Spec(\Fq)^{\Weil},
\]
其底层拓扑空间现在纤维化于 $S^1$ 之上，所有纤维都给出 $X_{\overline{\F}_q}$；绕圆周走一圈，就作用以 $\overline{\F}_q$ 的 Frobenius 自同构。更一般地，对 $\overline{\F}_q$ 上任一完满概形 $S$，可类似地定义
\[
X^{\Weil}_S=(X_S\times\R_{>0})/(\Frob_S\times q)^{\Z},
\]
它几乎、但并非完全是 $X^{\Weil}$ 的基变换（因为 $X^{\Weil}_S$ 事实上并不映射到 $S$）。我们还需要 Frobenius 改作用在 $X$ 上的版本，
\[
X^W=(X_{\overline{\F}_q}\times\R_{>0})/(\Frob_X\times q^{-1})^{\Z},
\]
（这里我们隐含地过渡到 $X$ 的完满化）。注意 $\pi_1(X^W)=W$ 就是 $X$ 的 Weil 群。
\end{remark}

\begin{definition}
带有 $I$ 条 leg 的 shtuka 模空间，是模空间
\[
f^I\colon\Sht^I\to(X^W)^I
\]
\srcpage{30}它把 $\overline{\F}_q$ 上任一完满概形 $S$ 送到 $I$ 个映射 $x_i\colon S\to X^W$，以及一个 $G$-丛\footnote{有一个稍显混乱、但本质上是自动的翻译，把 $x_i$ 的图像变成 $X^{\Weil}_S$ 的闭子集。}，后者定义在 $X^{\Weil}_S{\setminus}\bigcup_i\{\Gamma_{x_i}\}$ 上。
\end{definition}

因此，这参数化 $I$ 重穿孔三维流形上的 $G$-丛。这里的重要之处在于：沿 $\R_{>0}$-方向，这些 $G$-丛必须是常值的 —— 对任一概形 $S$，$S\times[0,1]$ 上的向量丛等价于 $S$ 上的向量丛。然而，当你越过一个穿孔时，可能发生一次修改。把上述 shtuka 的定义与更经典的定义联系起来，是一件令人愉快的练习。上述定义能够把关于修改种类的各种 shtuka 变体，合并为单个对象。\footnote{提示：对 $S=\Spec(\Fq)$，可以通过观察 $S^1$ 上的纤维来理解 $X^{\Weil}{\setminus}\{x_i\}$ 上的向量丛。在大多数纤维上，这不过是整个 $X_{\overline{\F}_q}$ 上的向量丛。把 $S^1$ 上的点稍作变动，由上述同伦不变性，该向量丛保持常值。在某些纤维上，有些点缺失；在那些纤维上，可以在缺失点处发生修改。最后，若绕圆周走一圈，最终回到原来的向量丛，但被 Frobenius 扭了一下。}

对任一 $V\in\Rep(\Gd^I)$，几何 Satake 等价给出 $\Sht^I$ 上的层 $\Sat_V$，它由某个适当的局部 Hecke 叠拉回而来。由 Xue（2025）的一个定理，shtuka 模空间的上同调层是光滑的。此外，$\pi_1(X^W)=W$ 是 Weil 群，且 $X^W$ 是固有的，故关于 $\pi_1$ 的 Künneth 定理成立，且 $\pi_1((X^W)^I)=W^I$。因此，我们得到
\[
f^I_!\Sat_V\in D(\Ql)^{BW^I}\subset D((X^W)^I,\Ql),
\]
从而得到所要的函子
\[
\Rep(\Gd^I)\to D(\Ql)^{BW^I}.
\]

我期望存在 $\Spec(\Z)$ 上的 shtuka 理论 —— 建立在把 shtuka 视为穿孔三维流形上的向量丛这一模板上 —— 用它最终能在数域情形执行类似的论证。

\begin{remark}
在 Betti 情形，定理 5.1 有一个适用于整个 $D(\Bun_G)$ 的类似物。即：沿用几何 Langlands 纲领中使用的记号，\footnote{除了把 $\Rep(\Gd)$ 换成 $D_{\qc}(\ast/\Gd)$ 之外；我们用 $\Rep(\Gd)$ 表示那个小 Abel 范畴。「Ran」这个名字在这里指「Ran 空间」，即「$X$ 的一切有限子集构成的空间」。我通常避免使用 Ran 空间，因为它是性态有些不良的几何对象，这里只用该记号与文献衔接。}令
\[
D_{\qc}(\ast/\Gd)_{\Ran}
\]
表示万有的 $C\in\operatorname{CAlg}(\PrL_{\e})$，它配有对称幺半函子
\[
\Rep(\Gd)\to C\otimes D(X).
\]
\srcpage{31}那么 Hecke 作用总能产生 $D_{\qc}(\ast/\Gd)_{\Ran}$ 在 $D(\Bun_G)$ 上的一个作用。在 Scholze（2026）与 Stefanich 的工作所发展的 Gestalten 框架中，$C$ 对应映射 Gestalt
\[
\Map(\mathrm{Gest}(D(X)),\ast/\Gd),
\]
即从 $D(X)$ 对应的 Gestalt 出发、指向 $\Gd$ 的分类叠的映射。换言之，这就是「Betti $\Gd$-局部系统的 Gestalt」的自然概念，因此我们将记
\[
\LocSys^{\Betti,\mathrm{Gest}}_{\Gd}=\Map(\mathrm{Gest}(D(X)),\ast/\Gd)=\mathrm{Gest}(D_{\qc}(\ast/\Gd)_{\Ran}).
\]
于是整个 $D(\Bun_G)$ 承载 $D_{\qc}(\ast/\Gd)_{\Ran}$ 的一个作用。隐含地把代数叠与其相联系的 Gestalt 等同起来，有自然的单射
\[
\LocSys^{\Betti}_{\Gd}\hookrightarrow\LocSys^{\Betti,\mathrm{Gest}}_{\Gd}.
\]
遗憾的是，在几何 Langlands 下，$D_{\Nilp}(\Bun_G)$ 会本质等价于 $D_{\qc}(\LocSys_{\Gd})$；但对整个 $D(\Bun_G)$ 而言并非如此。它并不等价于 $D_{\qc}(\LocSys^{\Betti,\mathrm{Gest}}_{\Gd})=D_{\qc}(\ast/\Gd)_{\Ran}$。
\end{remark}

\begin{remark}
在 étale 情形同样有 $D_{\qc}(\ast/\Gd)_{\Ran}$ 在整个 $D(\Bun_G)$ 上的作用的类似物。最干净的陈述，是把 Gestalten 的形式化更认真地使用，这次一直走到 $(\infty,2)$-范畴。即：给定的六函子形式化 $D$ 给出一个 transmutation 函子
\[
\Sch^{\mathrm{ft}}_k\to\mathrm{Gest}/\mathrm{Gest}(D(e))\colon S\mapsto[S]_D,
\]
它把任一 $S$ 送到 $S$ 上核的可表现 $(\infty,2)$-范畴的 Gestalt，参见 Scholze（2026，第 9 讲）。由下降，这延拓到叠。于是还可以考虑映射 Gestalt
\[
\LocSys^{\et,\mathrm{Gest}}_{\Gd}=\Map_{[\ast]_D}([X]_D,\ast/\Gd).
\]
它在 $[\ast]_D$ 上是 $1$-affine 的，从而对应 $\Spec(k)$ 上核的可表现 $(\infty,2)$-范畴 $K_D$ 中的一个交换代数。后者又可以描述为余端；它或许可称为 $D_{\qc}(\ast/\Gd)_{\Ran}$。于是 $D(\Bun_G)$ 成为 Gestalt $\LocSys^{\et,\mathrm{Gest}}_{\Gd}$ 上的范畴层，即 $K_D$ 中 $D_{\qc}(\ast/\Gd)_{\Ran}$ 上的模。\footnote{这类推理即使在 motivic 六函子形式化中也行得通。}
\end{remark}

我们稍后将用到局部系统的 Gestalt 的如下性质。

\begin{proposition}
$\LocSys^{\et,\mathrm{Gest}}_{\Gd}\to[\ast]_D$ 的对角是 $0$-affine 的，而自然映射
\[
\LocSys^{\et}_{\Gd}\times_{\mathrm{Gest}(D(e))}[\ast]_D\to\LocSys^{\et,\mathrm{Gest}}_{\Gd}
\]
是 $1$-étale 且 $1$-affine 的单射。\footnote{注意：$\LocSys^{\et}_{\Gd}$ 这个受限的局部系统叠是非常大的连通分支的无交并，特别地不是可数个仿射概形的并。为此，应取 Scholze（2026）中所用的截断基数 $\kappa$ 大于连续统，好让 $1$-étale 性的断言在技术上成立。}
\end{proposition}

\begin{proof}
\srcpage{32}首先，我们证明对角是 $0$-affine 的。由于 $[X]_D\to[\ast]_D$ 是 $0$-suave 的（因为 $X$ 上同调光滑），而 $\ast/\Gd$ 的对角是 $0$-affine 的，只需证明下述更一般的断言：在任一基 Gestalt $B$ 之上，若有一个在 $B$ 上 $0$-suave 的 Gestalt $S$，且 $T'\to T$ 是 $0$-affine 的，则 $\Map_B(S,T')\to\Map_B(S,T)$ 是 $0$-affine 的。为此，固定一个映射 $f\colon S\to T$，并考虑这样的函子：它把 $B$ 上的任一 $B'$ 送到 $f|_{B'}\colon S\times_B B'\to T$ 在 $T'\to T$ 上的提升构成的 anima。由于 $T'\to T$ 是 $0$-affine 的，这就是 $\operatorname{CAlg}(\cO(T)_1)$ 中映射
\[
(\cO(T')/\cO(T))_0\to(\cO(S\times_B B')/\cO(T))_0
\]
的 anima。但 $(\cO(S\times_B B')/\cO(T))_0$ 是 $(\cO(S\times_B B')/\cO(S))_0$ 在遗忘函子 $\operatorname{CAlg}(\cO(S)_1)\to\operatorname{CAlg}(\cO(T)_1)$（即沿 $f$ 的前推）下的像。此外，由 suave 基变换，$(\cO(S\times_B B')/\cO(S))_0\in\operatorname{CAlg}(\cO(S)_1)$ 是 $(\cO(B')/\cO(B))_0\in\operatorname{CAlg}(\cO(B)_1)$ 沿 $0$-suave 映射 $S\to B$ 的拉回。因此可以看到，$(\cO(S\times_B B')/\cO(T))_0$ 只依赖于 $(\cO(B')/\cO(B))_0\in\operatorname{CAlg}(\cO(B)_1)$。从而 $B'$ 只通过其 $0$-affine 截断进入。于是，该模空间在 $B$ 上是 $0$-affine 的。

现在容易看出映射
\[
\LocSys^{\et}_{\Gd}\times_{\mathrm{Gest}(D(e))}[\ast]_D\to\LocSys^{\et,\mathrm{Gest}}_{\Gd}
\]
是单射。事实上，这意味着它的对角是同构；但该对角是 $0$-affine 的，而在此情形可以比较模空间的描述。由于源与目标都是 $1$-affine 的，该映射是 $1$-affine 的。又，源是代数叠，从而是 $1$-étale 的；而目标的对角是 $0$-affine 的，从而是 $1$-étale 的；合起来便推出该映射是 $1$-étale 的。
\end{proof}

在 D-模的情形，情况略有不同。一般形式化仍把 $D(\Bun_G)$ 变成
\[
\LocSys^{\dR,\mathrm{Gest}}_{\Gd}=\Map(X_{\dR},\ast/\Gd)=\mathrm{Gest}(D_{\qc}(\ast/\Gd)_{\Ran})
\]
上的范畴层，其中 $X_{\dR}=X/(X\subset X\times X)^{\wedge}$ 是 $X$ 的 Simpson de Rham 叠。同样有单射
\[
\LocSys^{\dR}_{\Gd}\hookrightarrow\LocSys^{\dR,\mathrm{Gest}}_{\Gd}.
\]
然而，这一次无需对 $D(\Bun_G)$ 作任何修改：

\begin{theorem}[Gaitsgory，2015；Gaitsgory 与 Raskin，2025b]
$\LocSys^{\dR,\mathrm{Gest}}_{\Gd}$ 上的范畴层 $D(\Bun_G)$ 支撑在子空间
\[
\LocSys^{\dR}_{\Gd}\hookrightarrow\LocSys^{\dR,\mathrm{Gest}}_{\Gd}
\]
上。等价地，$D_{\qc}(\ast/\Gd)_{\Ran}$ 在 $D(\Bun_G)$ 上的作用经过 Verdier 商
\[
D_{\qc}(\ast/\Gd)_{\Ran}\to D_{\qc}(\LocSys^{\dR}_{\Gd})
\]
分解。
\end{theorem}

\begin{remark}
与 Betti 或 étale 情形不同，容许 $D_{\qc}(\LocSys^{\dR}_{\Gd})$ 的一个作用，并不强制 Hecke 作用沿 $X^I$-方向局部常值。归根结底，这是因为函子
\[
\Rep(\Gd)\to D_{\qc}(\LocSys^{\dR}_{\Gd})\otimes D_{\qc}(X_{\dR})
\]
\srcpage{33}（它编码万有的 de Rham $\Gd$-局部系统）并不经过
\[
D_{\qc}(\LocSys^{\dR}_{\Gd})\otimes\mathrm{IndLisse}(X_{\dR})\subset D_{\qc}(\LocSys^{\dR}_{\Gd})\otimes D_{\qc}(X_{\dR})
\]
分解。而这又是因为：把 $k$-代数 $A$ 送到 $D(A)\otimes D_{\qc}(X_{\dR})$ 中 lisse 对象的函子，不与 $A$ 中的基变换交换。
\end{remark}

这个定理极其非平凡。如同 D-模情形中关于几何 Langlands 的许多定理，它用到了这一框架的一些特殊之处。值得注意的是，有一个归纳地生成 $D(\Bun_G)$ 的对象的办法，它经由两种程序：其一是熟悉的 Eisenstein 级数构造，其二是 Kac--Moody 局部化。这给出产生尖点对象的一般办法。在 Langlands 纲领中，这是实际产生尖点对象最通用的途径。归根结底，它是几何 Langlands 对应中困难的「Galois 到自守」方向的来源。就手头的定理而言，人们证明这两个来源生成整个 $D(\Bun_G)$，且两个来源都支撑在
\[
\LocSys^{\dR}_{\Gd}\hookrightarrow\LocSys^{\dR,\mathrm{Gest}}_{\Gd}
\]
上。

\begin{remark}
虽然（层范畴意义下的）$\LocSys^{\dR,\mathrm{Gest}}_{\Gd}$ 有干净的生成元与关系描述，但尚不知道如何通过施加更多关系来给出子空间 $\LocSys^{\dR}_{\Gd}$ 的描述。这使得直接检验 $D(\Bun_G)$ 是否支撑在该子空间上变得困难。
\end{remark}

\section{$G=\Gm$}\srcpage{33}

本节把一切都对 $G=\Gm$ 的情形显式化。同样，我们假定 $X$ 的亏格 $g$ 满足 $g\ge 1$。为简单起见，我们还将固定一个基点 $x\in X$。此时 $\Pic=\Bun_{\Gm}$ 由
\[ \Pic=\Z\times J\times */\Gm \]
给出，其中 $J=\Pic^{0}$ 是 $X$ 的 Jacobian，一个主极化阿贝尔簇。因子 $*/\Gm$ 来自考虑 Picard 叠——若像讨论 Picard 概形时惯常所做的那样，在 $x\in X$ 处把线丛平凡化，我们得到的就会是 $\Z\times J$。具有幂零奇异支撑的条件，就简单地变成具有零奇异支撑的条件，而后者等价于该层是局部常值的。

因此，在 Betti 情形，我们得到
\[ D_{\Nilp}(\Pic)=\prod_{\Z}\Ind\Lisse(J)\otimes D(*/\C^{\times})=\prod_{\Z}D(*/H_{1}(X))\otimes D(*/\C^{\times}) \]
因为 $\pi_{1}(J)=H_{1}(X)=H_{1}(X,\Z)\cong\Z^{2g}$ 是 $X$ 的一阶同调群（$\Z$ 系数）。

\srcpage{34}在对偶一侧，我们计算 $\LocSys^{\Betti}_{\Gm}$。这是把任意 $k$-代数 $A$ 送到（作为群的）映射 $\pi_{1}\to A^{\times}$ 的叠，模去 $A^{\times}$ 经由共轭的平凡作用。正如我们在第 4 节对一般群所讨论的，
\[ \Hom(\pi_{1},\Gm)=\Gm^{2g}\times^{L}_{\Gm}*, \]
其中纤维积在导出意义下取。实际上，由于 $\Gm$ 交换，映射 $\Gm^{2g}\to\Gm$ 恒等于 $1$，于是
\[ \Hom(\pi_{1},\Gm)=\Gm^{2g}\times(*\times^{L}_{\Gm}*), \]
其中 $*\times^{L}_{\Gm}*=\Spec(\e[\epsilon])$ 是相应于 $\e[\epsilon]$ 的仿射导出概形，此处 $\epsilon$ 的次数为 $1$（因而平方为零，所以 $\e[\epsilon]=\e\oplus\e[1]$ 是平方零扩张）。总之，
\[ \LocSys^{\Betti}_{\Gm}=\Gm^{2g}\times(*\times^{L}_{\Gm}*)\times*/\Gm. \]

我们现在可以手算看出存在等价
\[ D_{\Nilp}(\Pic)\cong D_{\qc}(\LocSys^{\Betti}_{\Gm}). \]
事实上，这是三个等价之张量积：

\begin{enumerate}
\item 第一，
\[ \prod_{\Z}D(*)=D_{\qc}(*/\Gm): \]
这来自按权 $n\in\Z$ 对 $\Gm$-表示的分解。

\item 第二，
\[ D(*/H_{1}(X))\cong D_{\qc}(\Gm^{2g}): \]
事实上，$H_{1}(X)\cong\Z^{2g}$ 的一个表示是一个带有 $2g$ 个交换自同态的向量空间，它同时也是 $\Gm^{2g}$ 上的拟凝聚层。

\item 第三，
\[ D(*/\C^{\times})\cong D_{\qc}(\Spec(\e[\epsilon])). \]
事实上，$D(*/\C^{\times})$ 由单位对象沿光滑映射 $*\to*/\C^{\times}$ 的 $!$-推前紧生成，而该对象的自同态可算出是 $\C^{\times}$ 的同调，即 $\e[\epsilon]$。
\end{enumerate}

于是，在模式匹配的层面上，就得到所要的等价。我们注意到，对第三部分，我们需要 $\Spec(\e[\epsilon])$ 这一项，它要求人在导出代数几何的框架中工作。稍加仔细地做，可以证明谱作用把 $D_{\Nilp}(\Pic)$ 变成 $D_{\qc}(\LocSys^{\Betti}_{\Gm})$ 上的自由秩 $1$ 模。

\begin{remark}
上述所有等价都落入「Cartier 对偶」的一般模式，而在交换情形，Langlands 对偶确实可以看作 Cartier 对偶的一个应用。粗略地说，对任意类型的交换群概形或叠 $G$，都应有由（导出的）$\Hom(G,\Gm)$ 给出的 Cartier 对偶 $G^{*}$，此时应有等价 $D_{\qc}(*/G)\cong D_{\qc}(G^{*})$ 与 $D_{\qc}(*/G^{*})\cong D_{\qc}(G)$。例如，这出现在有限平坦群概形的经典情形中，\srcpage{35}但它对 $G=\Z$ 与 $G^{*}=\Gm$ 也成立。在阿贝尔簇 $A$ 的情形，对偶实际上是对偶阿贝尔簇 $A^{\vee}[-1]$，只是作了上同调次数 $1$ 的平移；此时该等价就变成 Fourier--Mukai 等价 $D_{\qc}(A)\cong D_{\qc}(A^{\vee})$。在 Gestalten 的框架中存在一般的 Cartier 对偶理论，参见 Scholze（2026）。
\end{remark}

在 étale 情形，故事以非常平行的方式展开。让我们更详细地讨论 de Rham 情形。在一侧，我们得到 $\Pic$ 上的 D-模范畴。同样，分解
\[ \Pic=\Z\times J\times */\Gm \]
把 $D(\Pic)$ 分解为三个因子 $D(\Z)$、$D(J)$ 与 $D(*/\Gm)$ 的张量积。因子 $D(\Z)$ 与 $D(*/\Gm)$ 等价于它们在 Betti 情形中的版本，并在 Langlands 对偶一侧以平行方式处理。然而 $D(J)=D_{\qc}(J_{\dR})$ 与它的 Betti 对应物颇不相同。

在 Langlands 对偶一侧，$\LocSys^{\dR}_{\Gm}$ 映射到 $\Pic$，而纤维由 $\Gamma(X,\OmX)$ 下的一个拟挠子给出。实际上，恰好命中 $\Pic$ 的次数 $0$ 部分 $J\times*/\Gm$。把一切展开，并同时记下 $H^{1}(X,\OmX)=k$，可以看出
\[ \LocSys^{\dR}_{\Gm}=EJ\times*/\Gm\times\Spec(k[\epsilon]), \]
其中
\[ 0\to V J\to EJ\to J\to 0 \]
是 Jacobian $J$ 的通用向量扩张，其核是对应于向量空间 $H^{0}(X,\OmX)$ 的向量群 $V J$。

现在结果表明，$J_{\dR}=J/(0\subset J)^{\wedge}$ 的 Cartier 对偶恰是 $EJ[-1]$——由于 $J$ 主极化，$J$ 本身对偶到 $J[-1]$，而形式交换群 $(0\subset J)^{\wedge}$ 对偶到 $V J$（因为一般地形式加法群是加法群的 Cartier 对偶）。因此 $D_{\qc}(J_{\dR})\cong D_{\qc}(EJ)$——这一 Cartier 对偶的实例归功于 Laumon（1996）。

注意 $\LocSys^{\Betti}_{\Gm}$ 中的因子 $\Gm^{2g}$ 被 $\LocSys^{\dR}_{\Gm}$ 中的因子 $EJ$ 替换。正如 Riemann--Hilbert 对应所预言的，它们有相同的 $\C$-点，甚至在 $\C$-点处有相同的形式完备化，但作为代数簇却非常不同。最重要的是，$\Gm^{2g}$ 是仿射的，因而有许多整体函数，而 $EJ$ 的整体截面只有 $k$。de Rham 情形中整体截面的这种稀少，实际上在几何 Langlands 定理证明的最后一步起了关键作用。

\section{Beilinson 谱投影子}\srcpage{35}

在 étale 理论中，Beilinson 谱投影子起关键作用。它给出一种从整个 $D(\Bun_G)$ 出发、以受控方式投影回性态良好的部分 $D_{\Nilp}(\Bun_G)$ 的办法。这有几个用途：\srcpage{36}

\begin{enumerate}
\item 它给出构造 $D_{\Nilp}(\Bun_G)$ 的对象的通用手段，特别地表明该范畴非平凡。值得注意的是，人们会把 Beilinson 谱投影子作用于 Whittaker 层，以构造 Langlands 函子。
\item 正如我们在下面所讨论的，它可用来产生 Hecke 特征层。
\item 计算 $D_{\Nilp}(\Bun_G)$ 上的 Frobenius 迹需要它。
\item 人们可以用它来证明定理 3.4 中的最终等价。目前 $D_{\Nilp}(\Bun_G)$ 指的是由定理 3.4 的性质 (1) 与 (2)（已知彼此等价）所定义的范畴。
\end{enumerate}

Arinkin、Gaitsgory、Kazhdan、Raskin、Rozenblyum 与 Varshavsky（2020b）中的原始构造，是一段关于范畴的复杂代数。我们将给出一个利用 Gestalt $\LocSys^{\et,\mathrm{Gest}}_{\Gd}$ 及其在整个 $D(\Bun_G)$ 上的作用的构造。这发生在基 Gestalt $[*]_D=\mathrm{Gest}(K_D)$ 之上：$K_D$ 是联系于六函子形式化 $D=\Ind D^{b}_{c}(-,\Ql)$（带 $(-_{*},-_{!})$-函子性）的对称幺半可表现 $(\infty,2)$-核范畴，$\mathrm{Gest}(K_D)$ 即其 Gestalt。

回忆注 5.7，我们有 $D(\Bun_G)\in K_D$，并带有 $\cO(\LocSys^{\et,\mathrm{Gest}}_{\Gd})_{1}=D_{\qc}(*/\Gd)_{\Ran}$ 的作用。此外，我们有子集
\[ i:\LocSys^{\et}_{\Gd}\times_{\mathrm{Gest}(D(\e))}[*]_D\subset\LocSys^{\et,\mathrm{Gest}}_{\Gd}. \]
我们于是可以把 $D(\Bun_G)$ 局部化到这个子集上——这是 $D(\Bun_G)$ 在 $D_{\qc}(\LocSys^{\et}_{\Gd})$ 作用于其上时的最佳逼近。因此，令
\[ D(\Bun_G)^{\mathrm{Hecke}}=D(\Bun_G)\otimes_{\cO(\LocSys^{\et,\mathrm{Gest}}_{\Gd})_{1}}\cO(\LocSys^{\et}_{\Gd})_{1}. \]

由于 $i$ 是 $1$-étale 且 $1$-proper（并且截断）的，它按 Scholze（2026，第 9 讲附录）容许 $!$-函子。因此，
\[ i^{*}:D(\Bun_G)\to D(\Bun_G)^{\mathrm{Hecke}} \]
容许一个截面 $i_{!}$，而 $D(\Bun_G)^{\mathrm{Hecke}}$ 是 $D(\Bun_G)$ 的收缩。另一方面，它在 $D(\Bun_G)$ 中的像必定包含于 Hecke-lisse 层之中，所以 $D(\Bun_G)^{\mathrm{Hecke}}$ 实际上是 $D_{\Nilp}(\Bun_G)\subset D(\Bun_G)$ 的收缩。\footnote{一个我们在此隐含使用的重要事实是：$D_{\Nilp}(\Bun_G)$ 是 restricted 的，即落在 $\PrL_{\e}\hookrightarrow K_D$ 的像中。} 由于后者已经是 $D_{\qc}(\LocSys^{\et}_{\Gd})$ 上的模，从而
\[ D(\Bun_G)^{\mathrm{Hecke}}=D_{\Nilp}(\Bun_G)\subset D(\Bun_G). \]

此外，函子 $i^{*}$ 现在定义 Beilinson 谱投影子
\[ \cP:D(\Bun_G)\to D_{\Nilp}(\Bun_G), \]
它是 $\cO(\LocSys^{\et,\mathrm{Gest}}_{\Gd})_{1}=D_{\qc}(*/\Gd)_{\Ran}$-线性的。

\begin{remark}
一个稍令人困惑之处是，Beilinson 谱投影子并不保证是包含
\[ i_{!}:D_{\Nilp}(\Bun_G)\subset D(\Bun_G) \]
的任何一侧的伴随。\srcpage{37}

人们应当想象 $i$ 既不是开浸入也不是闭浸入，而或许是局部闭浸入；在这种情形，$i_{!}$ 与 $i^{*}$ 不是伴随。
\end{remark}

作为紧生成范畴 $D(\Bun_G)$ 的收缩，形式上可知 $D_{\Nilp}(\Bun_G)$ 可对偶。实际上，可以证明它甚至是紧生成的：若使用与上面类似的手续，但把形式代数叠 $\LocSys^{\et}_{\Gd}$ 换成 Zariski 闭的代数子集，类似的映射 $i$ 将不再单射，但它将是 $0$-proper 的，因而函子 $i^{*}$ 有右伴随 $i_{*}$。这意味着在这一变体中 $i^{*}$ 保持紧对象，而这些紧对象给出 $D_{\Nilp}(\Bun_G)$ 的紧生成元。

关于对偶还可以说得更精确：在 $K_D$ 中，对象 $D(\Bun_G)$ 带有典范的自对偶数据，由余求值
\[ \Delta_{!}\Ql:\mathbf{1}\to D(\Bun_G)\otimes_{K_D}D(\Bun_G)=D(\Bun_G\times\Bun_G) \]
与求值
\[ \pi_{!}\Delta^{*}:D(\Bun_G)\otimes_{K_D}D(\Bun_G)\to\mathbf{1} \]
给出，即两者都由核 $\Delta_{!}\Ql$ 给出。这里 $\pi:\Bun_G\to*$ 表示投影，$\Delta$ 是对角。应用 Beilinson 谱投影子，可知 $D_{\Nilp}(\Bun_G)$ 也是自对偶的，相应的求值与余求值核由 Beilinson 谱投影子作用于 $\Bun_{G\times G}$ 上的 $\Delta_{!}\Ql$ 给出。\yizhu{原文此处作 $\Bun_{G\times G}$；按上下文，该核应是 $\Bun_G\times\Bun_G$ 上的 $\Delta_{!}\Ql$（前一式已如此写），此处照录原文。}{原文笔误（第 37 页；上下文推断）}

按 Gaitsgory 与 Raskin（2025a），包含
\[ D_{\Nilp}(\Bun_G)\subset D(\Bun_G) \]
保持紧对象；把它与 Arinkin、Gaitsgory、Kazhdan、Raskin、Rozenblyum 与 Varshavsky（2020b，第 4 节）结合起来——后者用 $(-^{!},-_{*})$-形式化中的伴随来刻画具有幂零奇异支撑的可构造层的性质——可以证明它实际上把紧对象送到 $(-_{*},-_{!})$-形式化中的 prim 层，从而给出一个等同
\[ D_{\Nilp}(\Bun_G)=\Ind D(\Bun_G)^{\prim}, \]
细节参见 Kuppel（2024）。我们注意到，在这一等同下，$D_{\Nilp}(\Bun_G)$ 的「神奇」自对偶不过是 prim 对偶——prim 对象经由取伴随，总带有反自等价。

借助 Beilinson 谱投影子，很容易定义 Hecke 特征层。即，对任意 $\Gd$-局部系统 $\varphi$，我们可以用 $D_{\qc}(\LocSys^{\et,\mathrm{Gest}}_{\Gd})$ 中由把 $i_{!}$ 作用于 $D_{\qc}(\LocSys^{\et}_{\Gd})$ 中的摩天大楼层而得到的对象来作用；这将投影到 $\varphi$ 对应的「Hecke 特征空间」上。遗憾的是，并不清楚这样能否得到非零对象。\srcpage{38}

\section{Frobenius 的迹}\srcpage{38}

取有限维向量空间 $V$ 的自同态 $f$ 的迹，这一通常操作有范畴的类比。更确切地说，更一般地，若 $V$ 是某个对称幺半范畴 $\cC$ 中的可对偶对象，且 $f\in\End(V)$，则可以定义
\[ \tr(f|V)\in\End_{\cC}(\mathbf{1}) \]
为复合
\[ \mathbf{1}\to V\otimes V^{*}\xrightarrow{f\otimes\id}V\otimes V^{*}\to\mathbf{1}. \]
这一构造有许多有趣的函子性性质，特别是迹的循环不变性，但我们在此不需要这些性质。

这也可以应用于 $\cC=\PrL_{\e}$ 以及任何可对偶的可表现 $\e$-线性 $(\infty,1)$-范畴。对可表现 $\e$-线性 $(\infty,1)$-范畴而言，可对偶这一条件原来是相当弱但非常有趣的条件；值得注意的是，Efimov（2025）最近把 $K$-理论推广到了这一框架。特别地，所有紧生成范畴都可对偶；并不要求同态空间具有某种有限性。\footnote{反过来，由于 $\PrL_{\e}$ 是 $(\infty,2)$-范畴，还有更强的 $2$-可对偶性概念。$2$-可对偶会推出紧对象之间同态空间的这类有限性条件。} 尽管存在有趣的可对偶非紧生成范畴的例子——例如流形上的层——我们实际上只需要紧生成的情形。

于是，给定任何可对偶的可表现 $\e$-线性 $(\infty,1)$-范畴 $\cC$ 与任何自函子 $F:\cC\to\cC$（在 $\PrL_{\e}$ 中，即保持余极限且 $\e$-线性），我们可以构造
\[ \tr(F|\cC)\in\End_{\PrL_{\e}}(\mathbf{1})=D(\e). \]
换言之，取范畴上的迹，得到的是一个向量空间（或者，由于我们完全在导出框架下工作，是向量空间导出范畴中的一个对象）。

\begin{example}
若 $F=\id$，这恢复 $\cC$ 的 Hochschild 同调。若对某个结合 $\e$-代数 $A$ 有 $\cC=D(A)$，则它由循环 Bar 复形
\[ [\dots\to A\otimes A\otimes A\to A\otimes A\to A] \]
计算，其中微分是 $a\otimes b\mapsto ab-ba$ 与 $a\otimes b\otimes c\mapsto ab\otimes c-a\otimes bc+ca\otimes b$，并对更多因子作适当推广。若 $\cC=D(A)$ 而 $F$ 是一般的，Morita 理论表明 $F$ 由某个 $A$-双模 $M$ 给出，而 $\tr(F|\cC)$ 由
\[ [\dots\to M\otimes A\otimes A\to M\otimes A\to M] \]
计算，其中微分是 $m\otimes a\mapsto ma-am$ 与 $m\otimes a\otimes b\mapsto ma\otimes b-m\otimes ab+ma\otimes b$\yizhu{上式末项原文亦作 $ma\otimes b$，与首项重复；对照前一条公式（末项为循环项 $ca\otimes b$），按循环 Bar 复形的通例此处应为 $mb\otimes a$。译文照录原文。}{原文笔误（第 38 页；上下文推断）}，利用双模结构，并同样作适当推广。
\end{example}

注意，若 $\cC$ 紧生成，则由 Schwede--Shipley 重构定理，它是形如 $D(A)$ 的范畴的递增并。由于迹构造\srcpage{39}与滤余极限交换，这在一般情形给出计算它的手段，至少当 $F$ 保持这一递增并的描述时如此。

可以在（导出的）向量空间 $\tr(F|\cC)$ 中描述一些显式的向量。即，若 $X\in\cC$ 是紧对象并带有一个映射 $\alpha:X\to F(X)$，就得到一个类
\[ [X,\alpha]\in\tr(F|\cC). \]

\begin{theorem}[Arinkin、Gaitsgory、Kazhdan、Raskin、Rozenblyum 与 Varshavsky，2021]
自然映射
\[ \tr(\Frob|D_{\Nilp}(\Bun_G))\to\cA_c(\Bun_G(\Fq)) \]
是同构。
\end{theorem}

利用 $(\infty,2)$-核范畴 $K_D$ 代替 $\PrL_{\e}$，并把迹形式化应用于其中的对象 $D(\Bun_G)$，实际上容易看出
\[ \tr(\Frob|D(\Bun_G))=\cA_c(\Bun_G(\Fq)). \]
事实上，在 $K_D$ 中，对象 $D(\Bun_G)$ 是自对偶的，其自对偶数据由 $\Delta_{!}\Ql\in D(\Bun_G\times\Bun_G)=D(\Bun_G)\otimes_{K_D}D(\Bun_G)$ 与
\[ (M,N)\mapsto\pi_{!}(M\otimes N):D(\Bun_G)\otimes_{K_D}D(\Bun_G)\to\mathbf{1} \]
给出。有了这些显式公式，容易算出 $\tr(\Frob|D(\Bun_G))$ 求值为 $\Bun_G^{\Frob=\id}=\Bun_G(\Fq)$ 的紧支撑上同调，即 $\cA_c(\Bun_G(\Fq))$。

由 $D_{\Nilp}(\Bun_G)$ 用 prim 对象所作的刻画，函子
\[ D_{\Nilp}(\Bun_G)\to D(\Bun_G) \]
在 $K_D$ 中有右伴随，从而诱导迹上的一个映射。这就是被证明为同构的映射
\[ \tr(\Frob|D_{\Nilp}(\Bun_G))\to\cA_c(\Bun_G(\Fq)). \]
由这一描述可知，它把带 Weil 下降数据 $\alpha$ 的紧对象 $\cF\in D_{\Nilp}(\Bun_G)$ 的类
\[ [\cF,\alpha]\in\tr(\Frob|D_{\Nilp}(\Bun_G)) \]
送到 $\cA_c(\Bun_G(\Fq))$ 中预期的函数
\[ x\mapsto\tr(\Frob_x|\cF_x). \]

虽然定理 8.2 是纯粹自守一侧的定理，它的证明却用到谱作用，特别是用到关于 shtuka 空间上同调的结果。

更确切地说，$D_{\Nilp}(\Bun_G)$ 的自对偶数据来自 $D(\Bun_G)$ 的自对偶数据，办法是应用 Beilinson 谱投影子。由于 Beilinson 谱\srcpage{40}投影子由谱作用、从而由 Hecke 作用得到，结果表明只需计算
\[ \tr(\Frob\circ H|D(\Bun_G\times X^{I})) \]
对 Hecke 算子 $H=H_V$，其中 $V\in\Rep(\Gd^{I})$。更确切地说，这必须在 $X^{I}$ 上的 $(\infty,2)$-核范畴中计算，因而结果不是（复形的）向量空间，而是 $D(X^{I})$ 的一个对象。Beilinson 谱投影子与从 $D(X^{I})$ 投影到 $\Ind\Lisse(X)^{\otimes I}$ 有关。所需的关键结果是
\[ \tr(\Frob\circ H|D(\Bun_G\times X^{I}))\in\Ind\Lisse(X)^{\otimes I}\subset D(X^{I}) \]
从而投影操作不起任何作用。

仿照 $H=\id$ 的论证的一个版本，一般地 $\tr(\Frob\circ H|D(\Bun_G\times X^{I}))$ 计算 shtuka 空间的紧支撑上同调。事实上，$H$ 来自 Hecke 对应
\[ \Hck^{I}_{G}\to\Bun_G\times\Bun_G\times X^{I}, \]
而人们必须与 Frobenius 的图像 $\Frob:\Bun_G\to\Bun_G$ 取纤维积。这一纤维积恰是一个 shtuka 空间，它参数化 $X$ 上的 $G$-丛连同它与自己的 Frobenius 拉回之间的一个同构，其中该同构允许在 $X$ 的 $I$ 个点处有极点。借助 Xue（2025）的定理——shtuka 空间的紧支撑上同调是 $X^{I}$ 上的 ind-lisse 层——即可得出结论，且它实际上落在 $\Ind\Lisse(X)^{\otimes I}\subset D(X^{I})$ 中。

\begin{remark}
对定理 1.7，还需要计算谱一侧的 Frobenius 迹。这是直接得多的计算，用到 Gaitsgory 与 Rozenblyum（2017）关于归纳凝聚层发展出的机器。
\end{remark}

\section{Whittaker 系数}\srcpage{40}

谱作用在各种设定下给出 $D_{\qc}(\LocSys_{\Gd})$ 在 $D(\Bun_G)$ 上的一个作用。在理想的情形下，这会使 $D(\Bun_G)$ 成为秩 $1$ 的自由模。尽管这一点会略有失效，仍然有必要造出所预期的基元素。换言之，我们寻找 $D(\Bun_G)$ 的一个特殊对象。这就是 Whittaker 层（亦称 Poincaré 对象），其定义如下。

回忆我们固定了 $\OmX$ 的一个平方根 $\Theta\in\Pic(X)$。特别地，把 Borel 子群 $B$ 用 $\rho(\OmX)=2\rho(\Theta)$ 共轭，我们得到 $B\subset G$ 的一个 $\Theta$-扭形式 $B^{\Theta}$。我们令 $U^{\Theta}\subset B^{\Theta}$ 为其单幂根。对 $G$ 的任一单根 $\alpha$，我们得到相应的投影
\[
\alpha\colon U^{\Theta}\to\OmX.
\]

\srcpage{41}令 $\psi\colon U^{\Theta}\to\OmX$ 为这些映射在全部单根上的和。过渡到分类空间并取从 $X$ 出发的映射叠，这诱导出一个映射
\[
\Bun^{\Theta}_{U}=\Bun_{U^{\Theta}}\to R\Gamma(X,\OmX)[1],
\]
其中我们滥用记号，把 $k$-向量空间与其相应的仿射空间等同。迹映射 $R\Gamma(X,\OmX)[1]\to k$ 由 Serre 对偶给出，它诱导出一个仍记作 $\psi$ 的映射
\[
\psi\colon \Bun^{\Theta}_{U}\to\A^{1}.
\]

在 $\A^{1}$ 上，通常有指数局部系统 $\exp\in D(\A^{1})$——在 D-模的情形，这是解为 $\exp(t)$ 的指数 D-模；在 étale 情形，这是 Artin--Schreier 层。在 Betti 情形它不存在，但可以添入，由此通向野 Betti 层的理论（Scholze，2025），并恢复 Tamarkin（2018）的一个构造。事实上，$\A^{1}$ 上这样的指数局部系统可以添入任何通常的层论，特别地可以添入动机，由此得到指数动机，参见 Fresán 与 Jossen（2025）以及 Gallauer 与 Pepin Lehalleur（2022）。因此，我们将不失一般性地假设我们的层论支持 $\exp$。

最后，利用对应
\[
\begin{tikzcd}
 & \Bun^{\Theta}_{U} \arrow[dl, "p"'] \arrow[dr, "\psi"] & \\
 \Bun_{G} & & \A^{1},
\end{tikzcd}
\]
我们得到 Whittaker 层
\[
\mathcal{W}=\mathcal{W}_{\psi}=p_{!}\psi^{*}\exp\in D(\Bun_{G}).
\]
这就是 $D(\Bun_{G})$ 中我们所期望的基对象。

\begin{remark}
先验地，可以取由单根诱导的映射的另一个线性组合 $\psi'\colon U^{\Theta}\to\OmX$。然而，只要系数非零，$\psi'$ 就会在 $T$ 的作用下与 $\psi$ 共轭，于是相应的 Whittaker 层也相同。归根结底，这种对 $\psi$ 的无关性还意味着：$\mathcal{W}_{\psi}$ 实际上不依赖层论中指数局部系统的选取，因此也可以在 Betti 设定下定义，而不必过渡到野 Betti 层。
\end{remark}

作为例子，取 $G=\mathrm{SL}_2$。此时 $\Bun^{\Theta}_{U}$ 分类扩张
\[
0\to\Theta\to E\to\Theta^{-1}\to 0;
\]
在同构意义下，这些扩张由其在 $\A^{1}$ 中的扩张类分类，这里再次用到 $H^{1}(X,\Theta^{2})=H^{1}(X,\OmX)=k$。每个扩张的自同构群是 $H^{0}(X,\OmX)$。

注意，若 $g\ge 2$，则 $\Theta$ 丰沛，$E$ 被 $\Theta$ 与 $\Theta^{-1}$ 给出的滤过被迫是 Harder--Narasimhan 滤过。事实上，对任意 $G$ 与 $g\ge 2$，映射 $p\colon\Bun^{\Theta}_{U}\to\Bun_{G}$\srcpage{42}分解为 $\Bun^{\Theta}_{U}\to\Bun^{\Theta}_{U}/T\to\Bun_{G}$，其中第二个映射是局部闭浸入，它分类如下轨迹：$G$-丛的 Harder--Narasimhan 滤过成为到 Borel 的约化，且相应的 $T$-丛同构于 $2\rho(\Theta)$。于是，对 $g\ge 2$，Whittaker 层极其显式。这一显式的知识实际上是几何 Langlands 对应证明最后一步的一个关键要素。

我们现在可以定义函子
\[
D_{\qc}(\LocSys_{\Gd})\to D_{\Nilp}(\Bun_G)\colon A\mapsto A\star\mathcal{P}(\mathcal{W}),
\]
其做法是：对 Whittaker 层施加 Beilinson 投影子，再使用谱作用。在 de Rham 情形，省去 Beilinson 谱投影子的施用，而使用整个 $D(\Bun_G)$。

\begin{definition}
粗 Langlands 函子是右伴随
\[
\mathbb{L}^{\coarse}_{G}\colon D_{\Nilp}(\Bun_G)\to D_{\qc}(\LocSys_{\Gd}).
\]
\end{definition}

由于谱作用保持紧对象，该右伴随仍保余极限。真正的 Langlands 函子应当是一个函子
\[
\mathbb{L}_{G}\colon D_{\Nilp}(\Bun_G)\to\IndCoh_{\Nilp}(\LocSys_{\Gd}),
\]
它映入凝聚复形的 Ind-范畴，这些复形在 Arinkin 与 Gaitsgory（2015）意义下的奇异支撑落在幂零锥中。\footnote{这与 $D_{\Nilp}(\Bun_G)$ 形式地类似，但两种奇异支撑的概念不同——分别指可构造层与凝聚层。}有一个遗忘函子
\[
\Psi\colon \IndCoh_{\Nilp}(\LocSys_{\Gd})\to D_{\qc}(\LocSys_{\Gd}),
\]
它是嵌入 $\Coh_{\Nilp}(\LocSys_{\Gd})\subset D_{\qc}(\LocSys_{\Gd})$ 的 Ind-扩张。函子 $\Psi$ 对自然的 $t$-结构是 $t$-正合的，并在最终余连通部分诱导同构——换言之，$\IndCoh_{\Nilp}$ 与 $D_{\qc}$ 之间的差别是 $\infty$-连通的。

\begin{theorem}[Gaitsgory 与 Raskin，2025b]
存在唯一的保余极限函子
\[
\mathbb{L}_{G}\colon D_{\Nilp}(\Bun_G)\to\IndCoh_{\Nilp}(\LocSys_{\Gd}),
\]
使得 $\mathbb{L}^{\coarse}_{G}=\Psi\circ\mathbb{L}_{G}$，并且它把紧对象送到最终余连通的对象。
\end{theorem}

\begin{remark}
让我们简要说明为何扩张是唯一的。由于所有函子都保余极限，它们完全由在紧对象上的行为决定。现在 Gaitsgory 与 Raskin（2025b）证明，$D(\Bun_G)$ 的紧对象在 $\mathbb{L}^{\coarse}_{G}$ 下的像最终余连通。由于 $\Psi$ 在最终余连通对象上是等价，我们就得到所期望的唯一提升。
\end{remark}

\srcpage{43}事实证明，人们可以在某种意义上显式地理解 Hecke 作用——从而谱作用——对 Whittaker 层做了什么。这用到几何 Casselman--Shalika 公式，它是几何 Satake 等价的一个近亲。回忆几何 Satake 等价只是阿贝尔范畴层面的等价；在导出层面，陈述变得更为微妙。反过来，几何 Casselman--Shalika 公式是一个完全的导出等价。

让我们陈述其局部版本。考虑 $LU^{\Theta}$ 在仿射 Grassmann 簇 $\Gr^{\Theta}_{G}$ 上的作用。这里 $\Gr^{\Theta}_{G}$ 是仿射 Grassmann 簇的 $\Theta$-扭形式 $LG^{\Theta}/L^{+}G^{\Theta}$；它参数化形式圆盘上的 $G$-挠子，连同与去心圆盘上的 $G$-丛 $2\rho(\Theta)$ 的一个等同。注意仍有一个映射
\[
\psi\colon LU^{\Theta}\to\A^{1},
\]
其做法是先对到单根空间的投影求和，得到到 $L\OmX$ 的映射，再用留数映射。在 $\A^{1}$ 上，如前有指数局部系统，于是我们考虑
\[
D^{!}(LU^{\Theta}\backslash{}^{\psi}\Gr^{\Theta}_{G})\subset D^{!}(\Gr^{\Theta}_{G}),
\]
即 $\Gr^{\Theta}_{G}$ 上在 $LU^{\Theta}$ 作用下 $\psi$-等变的 $!$-层的满子范畴。把 $\Gr^{\Theta}_{G}$ 写成有限维 Schubert 簇 $\Gr^{\Theta}_{G,\le\mu}$ 的余极限，就有
\[
D^{!}(\Gr^{\Theta}_{G})=\lim_{!}D(\Gr^{\Theta}_{G,\le\mu})=\operatorname{colim}_{!}D(\Gr^{\Theta}_{G,\le\mu}).
\]

\begin{remark}
$!$-层的出现或许可以看作一个迹象，表明人们必须使用 $(-_{!},-_{*})$-形式化。然而，$!$-层在 $(-_{*},-_{!})$-形式化中也有自然的意义。让我们先忘掉 $\Theta$-扭。我们有仿射 Grassmann 簇的自然的 transmutation
\[
[\Gr_{G}]_{D}=\operatorname{colim}_{\mu}[\Gr_{G,\le\mu}]_{D}.
\]
投影 $[\Gr_{G}]_{D}\to[*]_{D}$ 是 $1$-étale 的，但不是 $1$-proper 的。因此，在范畴层的层面，拉回有一个性质良好的左伴随，但没有性质良好的右伴随。所以，从 $[\Gr_{G}]_{D}$ 能构造出的好的层范畴，是把左伴随施用于 $[\Gr_{G}]_{D}$ 上的单位范畴层而得到的，这给出
\[
D^{!}(\Gr_{G})=\operatorname{colim}_{!}D(\Gr_{G,\le\mu}).
\]
类似地，$LU^{\Theta}$ 通过先 transmute 其有限维截断、再取逆向与正向极限，具有自然的 transmutation $[LU^{\Theta}]_{D}$。分类空间 $[*]_{D}/[LU^{\Theta}]_{D}$ 仍是 $1$-prim 的。\footnote{这归约到 $*/[L\A^{1}]_{D}$，在那里可以用 Cartier 对偶的一个应用来证明。}考察复合映射
\[
[LU^{\Theta}]_{D}\backslash[\Gr^{\Theta}_{G}]_{D}\xrightarrow{f}[LU^{\Theta}]_{D}\backslash[*]_{D}\xrightarrow{g}[*]_{D},
\]
我们可以定义
\[
D^{!}(LU^{\Theta}\backslash{}^{\psi}\Gr^{\Theta}_{G})=g_{1*}(f_{1\sharp}\mathbf{1}\otimes\psi),
\]
\srcpage{44}其中 $\psi$ 表示 $[LU^{\Theta}]_{D}\backslash[*]_{D}$ 上的可逆范畴层。它有自然的全忠实遗忘函子
\[
D^{!}(LU^{\Theta}\backslash{}^{\psi}\Gr^{\Theta}_{G})\to D^{!}(\Gr^{\Theta}_{G}),
\]
后者有左伴随。
\end{remark}

要给 $D^{!}(LU^{\Theta}\backslash{}^{\psi}\Gr^{\Theta}_{G})$ 一个干净的定义实际上很微妙，因为 $\Gr^{\Theta}_{G}$ 的归纳有限维本性与 $LU^{\Theta}$ 的归纳（无限维）本性混在一起。上一注的构造表明，Gestalten 的几何能够处理这种复杂性。

几何 Casselman--Shalika 公式给出 $D^{!}(LU^{\Theta}\backslash{}^{\psi}\Gr^{\Theta}_{G})$ 的一个描述。注意，从几何上看，Iwasawa 分解
\[
LG=LB\cdot L^{+}G=\bigsqcup_{\mu\in X_{*}(T)}LU\mu(t)L^{+}G
\]
表明商 $LU^{\Theta}\backslash\Gr^{\Theta}_{G}$ 具有按 $\mu\in X_{*}(T)$ 编号的层片的分层，层片由
\[
LU^{\Theta}\mu(t)L^{+}G^{\Theta}
\]
给出。并非所有层片都承载 $\psi$-等变的层；可以验证，这迫使 $\mu$ 支配。若 $\mu$ 支配，层片上的层等价于 $D(e)$。特别地，$\mu=0$ 给出基对象
\[
\Delta\in D^{!}(LU^{\Theta}\backslash{}^{\psi}\Gr^{\Theta}_{G}).
\]
另一方面，几何 Satake 范畴通过卷积作用在 $D^{!}(LU^{\Theta}\backslash{}^{\psi}\Gr^{\Theta}_{G})$ 上。

\begin{theorem}[几何 Casselman--Shalika]
作用
\[
D_{\qc}(*/\Gd)\to D^{!}(LU^{\Theta}\backslash{}^{\psi}\Gr^{\Theta}_{G})\colon V\mapsto\Sat(V)\star\Delta
\]
是范畴的等价。（事实上，在适当的解释下，得到 $[*]_{D}$ 上范畴层的一个同构。）
\end{theorem}

更一般地，对任何有限集 $I$，我们可以构造带 $LU^{\Theta}$-作用的 $\Gr^{\Theta}_{G,X^{I}}$，并把 $D^{!}(LU^{\Theta}\backslash{}^{\psi}\Gr^{\Theta}_{G,X^{I}})$ 定义为 $[X^{I}]_{D}$ 上的范畴层。在对角线之外，由同样的配方，它同构于 $D_{\qc}(*/\Gd^{I})$ 上的常范畴层；在对角线上，人们得到 $I$ 的商对应的同样情形。把所有函子性都处理妥当，就得到几何 Casselman--Shalika 公式的如下形式。

\srcpage{45}\begin{corollary}
设 $\Whit(G)$ 表示函子
\[
I\mapsto D^{!}(LU^{\Theta}\backslash{}^{\psi}\Gr^{\Theta}_{G,X^{I}})
\]
（视为 $[*]_{D}$ 上的范畴层）的余端。则 $D_{\qc}(*/\Gd)_{\Ran}$ 作用在 $\Whit(G)$ 上，且对 $\Delta$ 的作用诱导出等价
\[
\Whit(G)\cong D_{\qc}(*/\Gd)_{\Ran}.
\]
\end{corollary}

现在，对任何有限集 $I$，令 $\Bun^{U^{\Theta}}_{G,X^{I}}$ 表示 $G$-丛（等价地，经 $2\rho(\Theta)$ 作扭之后，即 $G^{\Theta}$-丛）的模空间，连带一个在 $I$ 点之外到 $U^{\Theta}$ 的约化。有对应
\[
\begin{tikzcd}
 & \Bun^{U^{\Theta}}_{G,X^{I}} \arrow[dl] \arrow[dr] & \\
 LU^{\Theta}\backslash\Gr^{\Theta}_{G,X^{I}} & & \Bun_{G}\times X^{I},
\end{tikzcd}
\]
或在 transmutation 之后类似的图表。此外，由 $\psi$ 给出的 $LU^{\Theta}\backslash*$ 上的范畴层，拉回到 $\Bun^{U^{\Theta}}_{G,X^{I}}$ 后变为平凡，因为在整体曲线上，留数之和为零。于是，由拉推我们得到一个函子\footnote{第二个函子是到 $D^{!}(\Bun_G)=D(\Bun_G)$ 的 $!$-推前。第一个函子用到：$\Bun^{U^{\Theta}}_{G,X^{I}}$ 是 $LU^{\Theta}\backslash\Gr^{\Theta}_{G,X^{I}}$ 沿映射 $*/L^{\mathrm{alg}}U^{\Theta}\to*/LU^{\Theta}$ 的拉回，其中 $L^{\mathrm{alg}}U^{\Theta}$ 把 $R$ 送到 $U^{\Theta}(X_{R}\setminus\bigcup_{i}\Gamma_{x_{i}})$；并且 $L^{\mathrm{alg}}U^{\Theta}$ 是 $0$-suave 的，于是 $*/L^{\mathrm{alg}}U^{\Theta}\to*$ 是 $0$-suave 且对角线 $0$-suave，从而在范畴层的层面，左、右伴随一致。}
\[
D^{!}(LU^{\Theta}\backslash{}^{\psi}\Gr^{\Theta}_{G,X^{I}})\to D^{!}(\Bun^{U^{\Theta}}_{G,X^{I}})\to D(\Bun_{G}).
\]
对全部 $I$ 取余端，我们得到函子
\[
\mathrm{Poinc}_{!}\colon \Whit(G)\to D(\Bun_{G}),
\]
它是 $D_{\qc}(*/\Gd)_{\Ran}$-线性的。由于源等价于 $D_{\qc}(*/\Gd)_{\Ran}$，这个函子等价于函子
\[
D_{\qc}(*/\Gd)_{\Ran}\to D(\Bun_{G})\colon A\mapsto A\star\mathcal{W}_{\psi},
\]
后者由对 Whittaker 层的谱作用给出；但现在它有一个纯粹在自守侧的描述！

\section{Eisenstein 级数}\srcpage{45}
\providecommand{\HH}{\mathbb{H}}% 本节新增宏，已在 reports/s10.md 登记

Langlands 理解自守形式空间的途径，基于一个归纳性的想法：先理解连续谱，\srcpage{46}而连续谱可以用 Levi 子群上的自守形式、借助 Eisenstein 级数来描述。回忆最经典的 Eisenstein 级数是如下求和
\[
E(\tau,s)=\frac{1}{2}\sum_{\substack{(m,n)\in\Z^{2}\\ \text{coprime}}}\frac{y^{s}}{|m+n\tau|^{2s}}
\]
其中 $\tau=x+iy$ 位于上半平面 $\HH=\{\tau\in\C\mid y=\operatorname{Im}(\tau)>0\}$ 内；该上半平面是 $\mathrm{SL}_2(\R)$ 经分式线性变换作用下的对称空间。而 $s\in\C$ 先验地要求 $\operatorname{Re}(s)>1$，以使该和收敛，但结果表明它可以亚纯延拓。由定义容易看出，$E(\tau,s)$ 下降为 $\HH/\mathrm{SL}_2(\Z)$ 上的函数。

Eisenstein 级数还有另一种写法。首先，
\[
E(\tau,s)=\sum_{\gamma\in\mathrm{SL}_2(\Z)/B(\Z)}\operatorname{Im}(\gamma(\tau))^{s},
\]
其中 $B(\Z)\subset\mathrm{SL}_2(\Z)$ 表示上三角矩阵构成的 Borel 子群（这里它就是 $\{\pm1\}\times\Z$）。事实上，稍作计算可知，对 $\tau=x+iy$ 有
\[
\operatorname{Im}(\gamma(\tau))=\operatorname{Im}\frac{a\tau+b}{c\tau+d}=\frac{y}{|c\tau+d|^{2}}.
\]

再作改写，我们可以考察由映射 $T\leftarrow B\to\mathrm{SL}_2$ 诱导的对应
\[
\begin{tikzcd}[column sep=large]
& B(\Z)\backslash B(\R)/\{\pm1\} \arrow[dl] \arrow[dr] & \\
T(\Z)\backslash T(\R)/\{\pm1\} & & \mathrm{SL}_2(\Z)\backslash\mathrm{SL}_2(\R)/\mathrm{SO}_2(\R)
\end{tikzcd}
\]
这里，右侧总取对极大紧子群的商。注意 $B(\R)/\{\pm1\}\to\mathrm{SL}_2(\R)/\mathrm{SO}_2(\R)$ 是同构。在此等同之下，投影 $B(\R)/\{\pm1\}\to T(\R)/\{\pm1\}=\R_{>0}$ 对应于投影
\[
\HH\to\R_{>0}\colon \tau=x+iy\mapsto\operatorname{Im}(\tau)=y.
\]
在这一解释下，我们从「$T=\Gm$ 上的自守形式」
\[
T(\Z)\backslash T(\R)/\{\pm1\}\cong\R_{>0}\to\C\colon y\mapsto|y|^{s}
\]
出发，把它拉回到 Borel 对应的空间 $B(\Z)\backslash B(\R)/\{\pm1\}$，再沿
\[
B(\Z)\backslash B(\R)/\{\pm1\}\to\mathrm{SL}_2(\Z)\backslash\mathrm{SL}_2(\R)/\mathrm{SO}_2(\R)
\]
的纤维求和，就得到 Eisenstein 级数 $E(\tau,s)$。

\srcpage{47}在这一解释下，显然可以看出如何把它推广得面目全非。用丛的模空间来说，我们可以考察对应
\[
\begin{tikzcd}[column sep=large]
& \Bun_B \arrow[dl,"p"'] \arrow[dr,"q"] & \\
\Bun_T & & \Bun_G,
\end{tikzcd}
\]
其中 $G$ 仍如前面，是带 Borel 子群 $B$ 与极大环面 $T$ 的分裂约化群。这诱导出函子
\[
\Eis_{!}=q_{!}p^{*}\colon D(\Bun_T)\to D(\Bun_G).
\]
实际上也可以把这个对应反向使用，从而得到常数项函子
\[
\CT_{!}=p_{!}q^{*}\colon D(\Bun_G)\to D(\Bun_T);
\]
在自守形式理论中，这对应于「沿尖点取 Fourier 展开中的常数项」。

更一般地，对 Levi 商为 $M$ 的抛物子群 $P$，我们得到对应
\[
\begin{tikzcd}[column sep=large]
& \Bun_P \arrow[dl,"p"'] \arrow[dr,"q"] & \\
\Bun_M & & \Bun_G.
\end{tikzcd}
\]
这给出类似的 Eisenstein 函子
\[
\Eis^{P}_{!}=q_{!}p^{*}\colon D(\Bun_M)\to D(\Bun_G)
\]
与常数项函子
\[
\CT^{P}_{!}=p_{!}q^{*}\colon D(\Bun_G)\to D(\Bun_M).
\]
显然，若把 $!$ 与 $*$ 分别换成 $*$ 与 $!$，可以得到这些函子彼此伴随的版本。有些出人意料的是，伴随性对上述这些函子本身也成立，只需把 $P$ 换成一个对立的抛物子群 $P^{-}$。按惯例这被称为「第二伴随」，因为它是 $p$ 进群表示论中 Bernstein 第二伴随的类比。\footnote{这一类比在 Hamann、Hansen 与 Scholze（2024）的设定中被处理得非常精确。}

\begin{theorem}[第二伴随；Drinfeld 与 Gaitsgory，2016]
函子 $\CT^{P^{-}}_{!}$ 是 $\Eis^{P}_{!}$ 的右伴随。该伴随关系在核的 $(\infty,2)$-范畴层面已经成立。
\end{theorem}

很容易给出该伴随的单位，即 $D(\Bun_M)$ 的自函子之间的一个映射
\[
\id\to\CT^{P^{-}}_{!}\circ\Eis^{P}_{!}.
\]
事实上，它由核的一个映射
\[
\Delta_{\Bun_M,!}\mathbf{1}\to f_{!}\mathbf{1}
\]
诱导。

\srcpage{48}这里 $f\colon\Bun_P\times_{\Bun_G}\Bun_{P^{-}}\to\Bun_M\times\Bun_M$ 是到两个 Levi 因子的投影。事实上，有开浸入
\[
*/M\hookrightarrow */P\times_{*/G}*/P^{-},
\]
它来自 $P\backslash G/P^{-}$ 中的开 Bruhat 胞腔 $*/M$；对 $X$ 取映射叠之后，这诱导开嵌入
\[
j\colon\Bun_M\to\Bun_P\times_{\Bun_G}\Bun_{P^{-}}.
\]
于是把 $f_{!}$ 作用于 $j_{!}\mathbf{1}\to\mathbf{1}$，即得所求的映射。

Drinfeld 与 Gaitsgory 在核的层面上给出该伴随余单位的一个几何构造；这用到一些与双曲局部化有关的想法，以及一族位于 $\A^{1}/\Gm$ 上的群 $\tilde G$，其一般纤维是 $G$，特殊纤维是 $P\times_M P^{-}$。\footnote{这一族 $\tilde G$ 也可以作为 Vinberg 幺半群的某一部分得到。}更确切地说，我们有一族
\[
\tilde G\to G\times G\times\A^{1}/\Gm,
\]
它诱导
\[
\Bun_{\tilde G}\to\Bun_G\times\Bun_G\times\A^{1}/\Gm,
\]
在对角嵌入与给出复合 $\Eis^{P}_{!}\circ\CT^{P^{-}}_{!}$ 的循环 $\Bun_P\times_{\Bun_M}\Bun_{P^{-}}$ 之间插值。借助一个余特殊化映射，由此可以造出所求的映射。

Eisenstein 级数的一个重要方面是：它们与 Hecke 算子大体上交换，从而把 Hecke 特征层送到 Hecke 特征层。严格说来，这并不完全成立。尽管如此，Campbell、Chen、Færgeman、Gaitsgory、Lin、Raskin 与 Rozenblyum（2024）在 D-模情形证明了下图差一个扭交换：
\[
\begin{tikzcd}[column sep=large,row sep=large]
D(\Bun_M) \arrow[r,"\mathbb{L}_M"] \arrow[d,"\Eis^{P}_{!}"'] & \IndCoh_{\Nilp}(\LocSys_{\Md}) \arrow[d,"\Eis^{P}_{\mathrm{spec}}"] \\
D(\Bun_G) \arrow[r,"\mathbb{L}_G"] & \IndCoh_{\Nilp}(\LocSys_{\Gd}).
\end{tikzcd}
\]
这里，谱 Eisenstein 函子指谱侧的一个类似构造。其证明用到：用 Whittaker 范畴 $\Whit(G)$ 对 Whittaker 层上的谱作用所作的描述，以及几何 Casselman--Shalika 公式与常数项的相容性。

\begin{remark}
Eisenstein 级数归根结底正是必须把 $D_{\qc}(\LocSys_{\Gd})$ 换成 $\IndCoh_{\Nilp}(\LocSys_{\Gd})$ 的原因。也就是说，函子 $\Eis^{P}_{!}$ 保持紧对象；但若把它的谱侧对应物看作 $D_{\qc}$ 上的函子，则后者不保持紧对象。真前推保持凝聚层的复形，却不保持完美复形。\srcpage{49}尽管如此，此时引入的奇异支撑方向必须沿给定抛物子群 $P$ 的单幂根而是幂零的。
\end{remark}

在上一个图中竖直地取右伴随，得到另一个图
\[
\begin{tikzcd}[column sep=large,row sep=large]
D(\Bun_M) \arrow[r,"\mathbb{L}_M"] & \IndCoh_{\Nilp}(\LocSys_{\Md}) \\
D(\Bun_G) \arrow[r,"\mathbb{L}_G"] \arrow[u,"\CT^{P^{-}}_{!}"] & \IndCoh_{\Nilp}(\LocSys_{\Gd}) \arrow[u,"\CT^{P^{-}}_{\mathrm{spec}}"']
\end{tikzcd}
\]
它先验地只到差一个自然归一化为止交换。Campbell、Chen、Færgeman、Gaitsgory、Lin、Raskin 与 Rozenblyum（2024）的一个主要定理是：这个图确实交换。其证明用到 Kac--Moody 局部化。

\section{Kac--Moody 局部化}\srcpage{49}
\providecommand{\Loc}{\mathrm{Loc}}% 本节新增宏，已在 reports/s11.md 登记
\providecommand{\Poinc}{\mathrm{Poinc}}% 本节新增宏，已在 reports/s11.md 登记

有了上一个图中那项重要的保留，我们迄今所讲的故事在任意层论中大体上都平行地成立；在很大的程度上，它在各种类似的设定中也成立，例如在 Fargues--Fontaine 曲线上。然而，几何 Langlands 猜想的证明用到一项只在 D-模理论中才具备的关键输入：Kac--Moody 局部化。这些想法可以追溯到 Beilinson 与 Drinfeld（1996）的开创性工作，他们从 oper 的概念出发，给出了 $\Bun_G$ 上 Hecke 特征层的一个构造。这是目前所知构造（尖点）Hecke 特征层最通用的办法，也是从 Galois 侧到自守侧这一方向的关键。这里，一个 oper 指一个（de Rham）$\Gd$-局部系统，连同一个把其底层的 $\Gd$-丛约化到 $\Ud^{\Th}$ 的约化，使得联络在相应的 $\Td$-丛上具有典范的形状。例如，当 $\Gd=\mathrm{SL}_2$ 时，它指的是秩 $2$ 的丛 $E$ 的一个滤过
\[
0\to\Th\to E\to\Th^{-1}\to 0,
\]
使得联络诱导出恒等映射
\[
\Th\to\Th^{-1}\otimes\Om^{1}=\Th.
\]

由于这一部分也是我最感生疏的部分，我在这里的讲述会比在本文其余部分更为印象式。基本想法是得到下图的一个版本
\[
\begin{tikzcd}[column sep=large,row sep=large]
D_{\qc}(*/\Gd)_{\Ran} \arrow[r,"\cong"] \arrow[d] & \Whit(G) \arrow[d] \\
D_{\qc}(\LocSys_{\Gd}) \arrow[r] & D(\Bun_G),
\end{tikzcd}
\]
只不过要交换 $G$ 与 $\Gd$ 的角色。

在谱侧，我们有某个隐约让人想起 $\Whit(G)$ 的东西：\srcpage{50}

\begin{definition}
形式圆盘 $\Spec(k[[t]])$ 上无单值 opers $\Op^{\monfree}_{\Gd}$ 的叠，是形式圆盘上带平坦联络的 $\Gd$-丛的模空间，并连带去心形式圆盘 $\Spec(k((t)))$ 上的一个 oper 结构。
\end{definition}

注意，形式圆盘上带平坦联络的 $\Gd$-丛全是平凡的；更确切地说，它们等价地由中心点处的一个 $\Gd$-挠子给出。叠 $\Op^{\monfree}_{\Gd}$ 无论在 Ind 意义下还是在 Pro 意义下都是极度无限型的。在所期望的「基本局部等价」的一侧，我们将有某个范畴
\[
\IndCoh^{*}(\Op^{\monfree}_{\Gd}),
\]
即 $\Op^{\monfree}_{\Gd}$ 上的归纳凝聚层构成的范畴；注意，由于这个对象极度无限维，要定义正确的归纳凝聚层概念并非易事。

在 Langlands 对应的另一侧，有一类由 Kazhdan--Lusztig 研究过的 Kac--Moody 表示构成的范畴，因而记作 $\KL(G)$。再次忽略各种 $\rho$-平移，它大致如下给出。回忆在有限维情形，李代数 $\gc$ 的表示也可以描述为 $*/\hat{G}$ 上的拟凝聚层，其中 $\hat{G}$ 是由 $\gc$ 生成的形式代数群。于是把表示从李代数 $\gc$ 积分到群 $G$，对应于沿 $*/\hat{G}\to*/G$ 延拓拟凝聚层。Kac--Moody 李代数的表示则与分类空间 $*/\LG^{\wedge}$ 上的拟凝聚层相关，这里 $\LG^{\wedge}$ 是 $\LG$ 在原点处的形式完备化。我们将关心其中能积分成正圈群 $\LpG\subset\LG$ 的表示的那些。这些表示也可以理解为分类空间
\[
*/(\LpG\subset\LG)^{\wedge}
\]
上的拟凝聚层，这里用的是 $\LG$ 沿闭子集 $\LpG\subset\LG$ 的形式完备化。

\begin{remark}
考虑映射
\[
f\colon */\LG\to */(\LG)_{\dR},
\]
即从 $*/\LG$ 到其 de Rham 叠的映射。于是 $f_{1\sharp}\mathbf{1}$ 定义了 $*/(\LG)_{\dR}$ 上的一个范畴层，其底范畴（即沿 $*\to*/(\LG)_{\dR}$ 的拉回）是 $*/\LG^{\wedge}$，也就是 Kac--Moody 表示。此外，若取 $f_{1\sharp}\mathbf{1}$ 的球向量，即在范畴层面取 $\LpG$-不变量，我们恰好得到
\[
D_{\qc}(*/(\LpG\subset\LG)^{\wedge}).
\]
\end{remark}

\begin{remark}
几何局部 Langlands 对应应当是 $(\infty,2)$-范畴之间的一个等价。在 D-模设定下，它把 $*/(\LG)_{\dR}$ 上的范畴层，与去心圆盘上 $\Gd$-局部系统叠上的（$\IndCoh_{\Nilp}$ 版的）拟凝聚范畴层联系起来。在此等价之下，上一个注中的表示应当对应于\srcpage{51}局部 opers 叠上的层范畴。而取球向量则对应于过渡到无单值 opers。
\end{remark}

在 Arinkin、Beraldo、Campbell、Chen、Færgeman、Gaitsgory、Lin、Raskin 与 Rozenblyum（2024）中，在 Feigin 与 Frenkel（1992）、Beilinson 与 Drinfeld（1996）以及许多其他人的工作基础上，证明了如下定理（见该文定理 0.1.2）：

\begin{theorem}[Arinkin、Beraldo、Campbell、Chen、Færgeman、Gaitsgory、Lin、Raskin 与 Rozenblyum，2024]
存在交换图
\[
\begin{tikzcd}[column sep=large,row sep=large]
D(\Bun_G)_{\crit} \arrow[r,"\mathbb{L}_G"] & \IndCoh_{\Nilp}(\LocSys_{\Gd}) \\
\KL(G)_{\crit} \arrow[r,"\mathrm{FLE}_G"] \arrow[u,"\Loc_G"] & \IndCoh^{*}(\Op^{\monfree}_{\Gd})_{\Ran} \arrow[u,"\Poinc^{\mathrm{spec}}_{\Gd}"']
\end{tikzcd}
\]
其中下方水平箭头是一个等价，即「基本局部等价」（在临界水平上）。
\end{theorem}

这里，下标 $\crit$ 指的是「临界水平」，它是通常的 $\rho$-平移的一个体现。就 $D(\Bun_G)$ 而言，由于我们对 $\Th$ 的选取，它其实并不改变范畴。

函子 $\Poinc^{\mathrm{spec}}_{\Gd}$ 有一个与我们此前见过的 $\Poinc_{!}$ 的定义非常相似的描述。也就是说，利用（一般的）整体 opers 的概念，即只在有限多个点之外才有 oper 结构，可以引入局部 opers 的空间与整体 $\Gd$-局部系统的叠之间的一个对应。然后施加通常的拉-推函子。

至于 $\Loc_G$ 的定义，我们注意到：由覆叠 $*/\LG\to*/(\LG)_{\dR}$ 得到的 $*/(\LG)_{\dR}$ 的表示，其整体类比是由覆叠 $\Bun_G\to\Bun_{G,\dR}$ 得到的 $D(\Bun_G)=D_{\qc}(\Bun_{G,\dR})$ 中的对象。注意，这只不过给出 $\Bun_G$ 上那个恰为微分算子代数本身的 D-模；而 Kac--Moody 局部化接着利用 $\LG$ 的李代数的表示论，切出更为具体的 D-模。

\begin{remark}
粗略地说，这个图使人能够构造有趣的 Hecke 特征层。也就是说，给定一个带 oper 结构的 $\Gd$-局部系统，就能在右下角造出一个对象。在基本局部等价之下，它可以被搬到自守侧，然后施加 Kac--Moody 局部化 $\Loc_G$，便得到 $\Bun_G$ 上一个有趣的 D-模，而它实际上将是给定 $\Gd$-局部系统的 Hecke 特征层。
\end{remark}
\srcpage{52}

\section{尖点形式}\srcpage{52}
\emergencystretch=3em

在 Arinkin、Beraldo、Campbell、Chen、Færgeman、Gaitsgory、Lin、Raskin 与 Rozenblyum（2024）以及 Campbell、Chen、Færgeman、Gaitsgory、Lin、Raskin 与 Rozenblyum（2024）两篇文章之后，几何 Langlands 等价在 Eisenstein 部分已经知道，尚待处理的是尖点部分。这里，尖点部分指常数项函子的核；但人们也知道，它同时也是由从 $\LocSys_{\Gd}$ 张量到不可约 $\Gd$-局部系统的开子空间 $\LocSys^{\irred}_{\Gd}$ 上而得到的部分。因此，我们考虑函子
\[
\begin{aligned}
\mathbb{L}^{\mathrm{cusp}}_G\colon{}& D(\Bun_G)\\
&{}\otimes_{D_{\qc}(\LocSys_{\Gd})}D_{\qc}(\LocSys^{\irred}_{\Gd})\\
&{}\to D_{\qc}(\LocSys^{\irred}_{\Gd}).
\end{aligned}
\]
我们只简要勾勒总体策略。

首先，证明 $\mathbb{L}^{\mathrm{cusp}}_G$ 既有左伴随 $(\mathbb{L}^{\mathrm{cusp}}_G)^L$，又有右伴随 $(\mathbb{L}^{\mathrm{cusp}}_G)^R$，且二者同构。更确切地说，由于 $\mathbb{L}^{\mathrm{cusp}}_G$ 的源与靶都是自对偶范畴，对偶函子 $(\mathbb{L}^{\mathrm{cusp}}_G)^{\vee}$ 可看作反方向的函子。随后证明：在相差同一个扭的意义下，这个函子既是 $\mathbb{L}^{\mathrm{cusp}}_G$ 的左伴随，也是它的右伴随。利用 $\mathbb{L}_G$ 与 Whittaker 系数的相容性，证明它是右伴随；利用与 Kac--Moody 局部化的相容性，证明它是左伴随。这正是 Arinkin、Beraldo、Chen、Færgeman、Gaitsgory、Lin、Raskin 与 Rozenblyum（2024）的内容。

最核心的输入是 Færgeman 与 Raskin（2025）的定理：$\mathbb{L}^{\mathrm{cusp}}_G$ 保守（用到「尖点形式自动缓增」）。他们的定理大致是这样证明的：把 Whittaker 系数等同于 $T^*\Bun_G$ 中正则幂零轨迹处的微茎（这一断言的精确版本参见 Nadler 与 Taylor（2025）），然后利用层的微局部性质，说明对缓增层该微茎的消失蕴涵完全消失。\footnote{对这一步骤，可以归约到 $D(\Bun_G)$ 中具有受限变差的部分，从而归约到具幂零奇异支撑的层。若正则幂零轨迹处的微茎消失，则奇异支撑受到更强的约束，而结果表明这与缓增性矛盾。}

这意味着，只需理解复合 $\mathbb{L}^{\mathrm{cusp}}_G\circ(\mathbb{L}^{\mathrm{cusp}}_G)^R$，并证明它同构于恒等函子。由于 $\mathbb{L}^{\mathrm{cusp}}_G$ 是 $D_{\qc}(\LocSys^{\irred}_{\Gd})$-线性的，且有左伴随，可知该函子由与 $\LocSys^{\irred}_{\Gd}$ 上某个完美复形 $A_{\Gd}$ 作张量给出。此外，由与 Kac--Moody 局部化、进而与 opers 的联系，$A_{\Gd}$ 由 oper 结构空间的同调给出。这已经证明 $A_{\Gd}$ 连通，且其零次同调群以一般 oper 结构空间的连通分支为自由生成元。另一方面，由于 $\mathbb{L}^{\mathrm{cusp}}_G$ 的左、右伴随相同，可知 $A_{\Gd}$ 自对偶。又因它连通，实际上它是 0 次的一个向量丛。通过「以一般 oper 结构空间的连通分支为自由生成元」这一解释，它实际上是带联络的向量丛，从而是 $\LocSys^{\irred}_{\Gd}$ 上的局部系统。\srcpage{53}

Beraldo、Kazhdan 与 Schlank（2022）的近期定理表明：对经典群，一般 oper 结构空间非空且连通，于是经典群的证明完成。

另一方面，至少在 $g\ge 2$、并排除 $g=2$ 且 $G$ 为 $A_1$ 型的情形时\footnote{其余情形的尖点部分可以容易地手工处理。}，可以计算 $\LocSys^{\irred}_{\Gd}$ 的基本群；它与 $G$ 的中心的连通分支群有关。若 $G$ 的中心连通，则看出 $A_{\Gd}$ 是 $\LocSys^{\irred}_{\Gd}$ 的结构层的若干拷贝的直和；注意后者（即 $\LocSys^{\irred}_{\Gd}$）实际上是连通的。\footnote{这一连通性是另一个关键要素，且与算术 Langlands 的情形大不相同：在几何 Langlands 纲领中，尖点形式及其对偶 $L$-参数同样构成一个「连续谱」。}此外，$\LocSys^{\irred}_{\Gd}\subset\LocSys_{\Gd}$ 的补集余维至少为 2，且全空间是局部完全交；于是关于深度（depth）的考量表明，只需计算 $\LocSys_{\Gd}$ 上某个 $H^0$ 的维数。但把伴随展开，这由 Whittaker 层的自同态给出。由于 Whittaker 层完全显式，其自同态极易计算，且为一维，正如所愿。若 $G$ 的中心不连通，本质上同样的论证可行，只需一些额外的扭。

\begin{remark}
这一最后论证的关键之处在于：在 de Rham 情形，局部系统的叠只有极少的整体截面，因此最终可以把证明「处处重数一」的问题，转化为证明「整体截面层面重数一」。只有后一陈述适合在自守侧做这一计算。
\end{remark}

\emergencystretch=0pt
\section*{参考文献}\srcpage{53}

\begin{list}{}{\setlength{\leftmargin}{2em}\setlength{\itemindent}{-2em}\setlength{\labelwidth}{0pt}\setlength{\labelsep}{0pt}\setlength{\itemsep}{0.3em}\setlength{\parsep}{0pt}\setlength{\topsep}{0.4em}}

\item Avraham Aizenbud and Nir Avni (2016). “Representation growth and rational singularities of the moduli space of local systems”, \textit{Invent. Math.} \textbf{204} (1), pp. 245–316.
\item Dima Arinkin, Dario Beraldo, Justin Campbell, Lin Chen, Joakim Færgeman, Dennis Gaitsgory, Kevin Lin, Sam Raskin, and Nick Rozenblyum (2024). \textit{Proof of the geometric Langlands conjecture II: Kac-Moody localization and the FLE}. arXiv: 2405.03648.
\item Dima Arinkin, Dario Beraldo, Lin Chen, Joakim Færgeman, Dennis Gaitsgory, Kevin Lin, Sam Raskin, and Nick Rozenblyum (2024). \textit{Proof of the geometric Langlands conjecture IV: ambidexterity}. arXiv: 2409.08670.
\item Dima Arinkin and Dennis Gaitsgory (2015). “Singular support of coherent sheaves and the geometric Langlands conjecture”, \textit{Selecta Math. (N.S.)} \textbf{21} (1), pp. 1–199.
\item Dima Arinkin, Dennis Gaitsgory, David Kazhdan, Sam Raskin, Nick Rozenblyum, and Yakov Varshavsky (2020a). \textit{Duality for automorphic sheaves with nilpotent singular support}. arXiv: 2012.07665.
\srcpage{54}
\item Dima Arinkin, Dennis Gaitsgory, David Kazhdan, Sam Raskin, Nick Rozenblyum, and Yakov Varshavsky (2020b). \textit{The stack of local systems with restricted variation and geometric Langlands theory with nilpotent singular support}. arXiv: 2010.01906.
\item ——— (2021). \textit{Automorphic functions as the trace of Frobenius}. arXiv: 2102.07906.
\item James Arthur (1989). “Unipotent automorphic representations: conjectures”, in: \textit{Orbites unipotentes et représentations, II}. Astérisque 171-172, pp. 13–71.
\item Alexander Beilinson (2016). “Constructible sheaves are holonomic”, \textit{Selecta Math. (N.S.)} \textbf{22} (4), pp. 1797–1819.
\item Alexander Beilinson and Vladimir Drinfeld (1996). “Quantization of Hitchin's fibration and Langlands' program”, in: \textit{Algebraic and geometric methods in mathematical physics (Kaciveli, 1993)}. Math. Phys. Stud. (Kluwer Acad. Publ.), vol. 19, pp. 3–7.
\item David Ben-Zvi (2026). “What is the geometric Langlands correspondence about?”, \textit{Current Events Bulletin of the AMS}, pp. 4–36.
\item David Ben-Zvi and David Nadler (2013). “Loop spaces and representations”, \textit{Duke Math. J.} \textbf{162} (9), pp. 1587–1619.
\item ——— (2018). “Betti geometric Langlands”, in: \textit{Algebraic geometry: Salt Lake City 2015}. Proc. Sympos. Pure Math. (AMS), vol. 97.2, pp. 3–41.
\item Dario Beraldo, David Kazhdan, and Tomer M. Schlank (2022). \textit{Contractibility of the space of generic opers for classical groups}. arXiv: 1801.00655v2.
\item Kevin Buzzard and Toby Gee (2014). “The conjectural connections between automorphic representations and Galois representations”, in: \textit{Automorphic forms and Galois representations. Vol. 1}. London Math. Soc. Lecture Note Ser. (Cambridge Univ. Press), vol. 414, pp. 135–187.
\item Justin Campbell, Lin Chen, Joakim Færgeman, Dennis Gaitsgory, Kevin Lin, Sam Raskin, and Nick Rozenblyum (2024). \textit{Proof of the geometric Langlands conjecture III: compatibility with parabolic induction}. arXiv: 2409.07051.
\item Vladimir Drinfeld (1983). “Two-dimensional $l$-adic representations of the fundamental group of a curve over a finite field and automorphic forms on $GL(2)$”, \textit{Amer. J. Math.} \textbf{105} (1), pp. 85–114.
\item Vladimir Drinfeld and Dennis Gaitsgory (2016). “Geometric constant term functor(s)”, \textit{Selecta Math. (N.S.)} \textbf{22} (4), pp. 1881–1951.
\item Alexander Efimov (2025). \textit{K-theory and localizing invariants of large categories}. arXiv: 2405.12169v3.
\item Matthew Emerton, Toby Gee, and Eugen Hellmann (2025). \textit{An introduction to the categorical p-adic Langlands program}. arXiv: 2210.01404v5.
\item Joakim Færgeman and Sam Raskin (2025). “Non-vanishing of geometric Whittaker coefficients for reductive groups”, \textit{J. Amer. Math. Soc.} \textbf{38} (4), pp. 919–995.
\item Laurent Fargues (2025). “Geometrization of the local Langlands correspondence: an overview”, \textit{J. Math. Sci. Univ. Tokyo} \textbf{32} (2), pp. 157–240.
\item Laurent Fargues and Peter Scholze (2024). \textit{Geometrization of the local Langlands correspondence}, to appear in Astérisque. arXiv: 2102.13459.
\srcpage{55}
\item Boris Feigin and Edward Frenkel (1992). “Affine Kac-Moody algebras at the critical level and Gelfand-Dikiĭ algebras”, in: \textit{Infinite analysis, Part A, B (Kyoto, 1991)}. Adv. Ser. Math. Phys. (World Sci. Publ.), vol. 16, pp. 197–215.
\item Javier Fresán and Peter Jossen (2025). “Exponential motives”. \url{https://javier.fresan.perso.math.cnrs.fr/expmot.pdf}.
\item Dennis Gaitsgory (2015). “Outline of the proof of the geometric Langlands conjecture for $GL_2$”, \textit{Astérisque} (370), pp. 1–112.
\item Dennis Gaitsgory and Sam Raskin (2025a). \textit{Geometric Langlands in positive characteristic from characteristic zero}. arXiv: 2508.02237.
\item ——— (2025b). \textit{Proof of the geometric Langlands conjecture I: construction of the functor}. arXiv: 2405.03599v3.
\item ——— (2026). \textit{Proof of the geometric Langlands conjecture V: the multiplicity one theorem}. arXiv: 2409.09856v3.
\item Dennis Gaitsgory and Nick Rozenblyum (2017). \textit{A study in derived algebraic geometry. Vol. I. Correspondences and duality}. Math. Surveys and Monographs (AMS), vol. 221.
\item Martin Gallauer and Simon Pepin Lehalleur (2022). “Exponentiation of coefficient systems and exponential motives”. arXiv: 2211.17247.
\item Linus Hamann, David Hansen, and Peter Scholze (2024). \textit{Geometric Eisenstein series I: finiteness theorems}. arXiv: 2409.07363.
\item David Hansen and Peter Scholze (2023). “Relative perversity”, \textit{Commun. Am. Math. Soc.} \textbf{3}, pp. 631–668.
\item Claudius Heyer and Lucas Mann (2024). \textit{6-Functor Formalisms and Smooth Representations}. arXiv: 2410.13038.
\item Tim Kuppel (2024). \textit{Proper Sheaves for Schemes and Stacks}. Master Thesis, University of Bonn.
\item Vincent Lafforgue (2018). “Chtoucas pour les groupes réductifs et paramétrisation de Langlands globale”, \textit{J. Amer. Math. Soc.} \textbf{31} (3), pp. 719–891.
\item Vincent Lafforgue and Xinwen Zhu (2019). \textit{Décomposition au-dessus des paramètres de Langlands elliptiques}. arXiv: 1811.07976v2.
\item Gérard Laumon (1987). “Correspondance de Langlands géométrique pour les corps de fonctions”, \textit{Duke Math. J.} \textbf{54} (2), pp. 309–359.
\item ——— (1996). \textit{Transformation de Fourier géneralisée}. arXiv: alg-geom/9603004.
\item Yifeng Liu and Weizhe Zheng (2024). \textit{Enhanced six operations and base change theorem for higher Artin stacks}. arXiv: 1211.5948.
\item Jacob Lurie (2017). “Higher Algebra”. \url{https://www.math.ias.edu/~lurie/papers/HA.pdf}.
\item Lucas Mann (2022). \textit{A p-Adic 6-Functor Formalism in Rigid-Analytic Geometry}. arXiv: 2206.02022.
\item Ivan Mirković and Kari Vilonen (2007). “Geometric Langlands duality and representations of algebraic groups over commutative rings”, \textit{Ann. of Math. (2)} \textbf{166} (1), pp. 95–143.
\srcpage{56}
\item David Nadler and Jeremy Taylor (2025). “The Whittaker functional is a shifted microstalk”, \textit{Transform. Groups} \textbf{30} (3), pp. 1425–1450.
\item David Nadler and Zhiwei Yun (2019). “Spectral action in Betti geometric Langlands”, \textit{Israel J. Math.} \textbf{232} (1), pp. 299–349.
\item Sam Raskin (2025). “Conjectures of Arthur and Ramanujan for unramified automorphic forms: announcement”. \url{https://www.samraskin.net/icm.pdf}.
\item Timo Richarz and Jakob Scholbach (2021). “The motivic Satake equivalence”, \textit{Math. Ann.} \textbf{380} (3-4), pp. 1595–1653.
\item Peter Scholze (2024). \textit{Geometrization of the real local Langlands correspondence}.
\item ——— (2025). “Wild Betti sheaves”, in: \textit{Algebraic Geometry and the Langlands Program, a volume in honour of Gérard Laumon}.
\item ——— (2026). “Geometry and Higher Category Theory”. \url{https://people.mpim-bonn.mpg.de/scholze/Gestalten.pdf}.
\item Dmitry Tamarkin (2018). “Microlocal condition for non-displaceability”, in: \textit{Algebraic and analytic microlocal analysis}. Springer Proc. Math. Stat. vol. 269, pp. 99–223.
\item Thibaud van den Hove (2024). \textit{The integral motivic Satake equivalence for ramified groups}. arXiv: 2404.15694.
\item André Weil (1974). \textit{Basic number theory}. Third edition. Die Grundlehren der mathematischen Wissenschaften (Springer), Bd. 144.
\item Cong Xue (2025). “Smoothness of cohomology sheaves of stacks of shtukas”, in: \textit{Algebraic Geometry and the Langlands Program, a volume in honour of Gérard Laumon}.
\item Xinwen Zhu (2024). “A note on integral Satake isomorphisms”, in: \textit{Arithmetic geometry}. Tata Inst. Fundam. Res. Stud. Math, vol. 41, pp. 469–489.
\item ——— (2025). “Coherent sheaves on the stack of Langlands parameters”, in: \textit{The Langlands program. II—geometrization of the Langlands correspondence}. Proc. Sympos. Pure Math. (AMS), vol. 112, pp. 39–123.


\end{list}

\bigskip
\begin{flushright}
Peter Scholze\\
Max-Planck-Institut für Mathematik\\
Vivatsgasse 7\\
53111 Bonn\\
\textit{E-mail} : scholze@mpim-bonn.mpg.de
\end{flushright}


\end{document}
